% !TeX root = ../../Book2.tex
\chapter{极大函数、弱型估计与插值}
\label{b2:ch:17}
\BookChapterNumber{b2:ch:17}
\begin{chapterabstract}
估计一个算子时，我们往往只能先掌握两个指数：一端来自绝对积分，另一端来自平方范数或本性上确界。中间的指数怎样得到？端点的积分发散以后，还能保留什么？本章从局部平均的极大函数出发，用分布函数记录大值出现的范围，再分别通过大小截断与复参数变化建立两种插值方法。弱空间、Lorentz 空间和 \(K\) 泛函说明了这些估计背后的函数空间结构。
\end{chapterabstract}

\begin{learninggoals}\begin{itemize}
\item\label{b2:item:17-goal-1}比较中心、非中心与二进平均，并完成弱型估计和微分论证。
\item\label{b2:item:17-goal-2}由分布函数计算强、弱与 Lorentz 空间的可积条件。
\item\label{b2:item:17-goal-3}选择大小截断并完成 Marcinkiewicz 的高度积分。
\item\label{b2:item:17-goal-4}用三线估计证明 Riesz--Thorin，并解释 Fourier 变换的延拓。
\end{itemize}\end{learninggoals}

单位球的示性函数有有限积分，但其极大平均在远处只按 \(|x|^{-d}\) 衰减。这个简单例子说明：一个自然算子可能在端点失去可积性，同时仍有精确的超水平集合控制。后面的奇异积分也呈现同样现象，只是核的正平均将换成正负抵消。

\ProseParagraphBreak

本章需要 Lebesgue 积分、Hölder 不等式以及 \(L^p\) 收敛的基本知识。遇到无限求和或换序时，我们会指出非负性、可积控制或有限测度截断在哪里使用。极大函数的微分论证还要用连续紧支集函数的 \(L^1\) 稠密性；Fourier 应用则使用\secref{b2:sec:0401}已经建立的 Plancherel 等距性。复插值所需的复分析知识限于全纯函数的局部性质与最大模原理，\secref{b2:sec:1704}先说明其使用方式，再证明条带上的三线引理。

\ProseParagraphBreak

阅读时可先沿“极大函数的弱端点---分布函数的高度积分---中间强型”这一条路线走完\secref{b2:sec:1701}至\secref{b2:sec:1703}，看清为什么保留输入的大值部分能够消除发散。随后学习\secref{b2:sec:1704}，比较线性结构带来的另一种插值。Lorentz 空间的第二指数和最后的 \(K\) 泛函，可以在掌握这两条证明之后再回读；极大函数及同指数插值的证明并不依赖这两个补充视角。\secref{b2:sec:1701}中的二进立方体分解将在\secref{b2:sec:1802}直接用于 Calderón--Zygmund 分解。

\ProseParagraphBreak

继续阅读可选 Elias M. Stein 的《Singular Integrals and Differentiability Properties of Functions》\cite{Stein1970Singular}，把极大函数与后续奇异积分联系起来；Loukas Grafakos 的《Classical Fourier Analysis》\cite{Grafakos2014Classical}则可用于进一步学习弱型、重排和插值。本章所需的具体估计与证明在正文展开。

% !TeX root = ../../Book2.tex
\section{Hardy–Littlewood 极大算子回顾}
\label{b2:sec:1701}

固定尺度的平均会平滑函数，却可能漏掉比这个尺度窄得多的尖峰。若要同时讨论许多尺度的逼近，就希望用同一个函数控制所有这些平均。极大算子的做法是：先在每一点计算各个尺度上的绝对平均，再取其中的上确界。这个上确界允许观察尺度随位置改变，因而不能简单视为一个固定核的卷积。本节固定欧氏空间的 Lebesgue 测度，\(d\geq1\) 为整数。

\subsection{第一卷结果回顾}
\label{b2:subsec:170101}

对实或复函数 \(f\in L^1_{\mathrm{loc}}(\mathbb R^d)\)，约定
\[
 Mf(x)\coloneqq\sup_{r>0}\dfrac1{|B(x,r)|}\int_{B(x,r)}|f(y)|\,\mathrm dy.
\]
\symidx{\(Mf\)：Hardy--Littlewood 中心极大函数}
这是本卷的\term[hardy-littlewood-jida-hanshu]{Hardy--Littlewood 极大函数}[Hardy--Littlewood maximal function]，相应映射 \(f\mapsto Mf\) 称为极大算子。先取绝对值，是为了不让输入自身的正负抵消遮住误差大小；再除以球体积，是为了使常函数的平均保持原高度。它一般不是线性算子，但满足
\[M(f+g)\leq Mf+Mg,\quad M(cf)=|c|\cdot Mf\quad(c\in\mathbb C).\]
这两条性质分别来自积分三角不等式和绝对齐次性，后面的实插值会利用这种结构。

\ProseParagraphBreak

还须确认无穷多个平均的上确界是可测的。记 \(A_rf(x)=|B(x,r)|^{-1}\int_{B(x,r)}|f|\)。固定 \(r>0\)，当中心从 \(x\) 移到 \(x+h\) 时，两个积分的差不超过
\[\int_{B(x+h,r)\,\triangle\,B(x,r)}|f|.\]
这些球位于一个固定有界区域内，对称差的体积随 \(h\to0\) 趋零；局部可积函数的积分对小测度集合绝对连续，故 \(A_rf\) 对中心连续。固定中心时，令半径趋于某个正数，两个球之间的薄环体积也趋零，且分母体积连续，因此 \(r\mapsto A_rf(x)\) 也连续。于是
\[Mf(x)=\sup_{r\in\mathbb Q,\ r>0}A_rf(x)\]
是可数个可测函数的上确界。另一方面，\(\{Mf>t\}\) 是各固定半径的开超水平集合之并，因而还是开集。局部可积性通过积分的绝对连续性，转化为了平均值对球心与半径的连续性。

\ProseParagraphBreak

一个能直接算完的例子是实直线上的
\[f=\mathbf1_{(0,1)}+\mathbf1_{(2,3)}.\]
在空隙中点 \(x=3/2\)，半径 \(r\leq1/2\) 的平均为零；若 \(1/2<r<3/2\)，球内与两段支集相交的总长度为 \(2r-1\)，故平均为 \(1-1/(2r)\)，随 \(r\) 增大；若 \(r\geq3/2\)，总质量固定为 \(2\)，平均变为 \(1/r\)，随 \(r\) 减小。因此 \(Mf(3/2)=2/3\)，尽管该点附近 \(f=0\)。极大平均会从较大邻域看见质量，不能把它与函数本身在该点的大小混同。

\ProseParagraphBreak

若某个平均超过高度 \(t\)，它所在的球必须含有至少 \(t\) 乘体积的质量。困难在于这些球会重叠，不能直接把它们的积分相加。下面先选出不交的一部分，再用固定倍数的扩大球补回覆盖。

\Needspace{6\baselineskip}
\begin{theorem}[Hardy--Littlewood]\label{b2:thm:17-maximal-weak}
设整数 \(d\geq1\)，\(f\in L^1(\mathbb R^d;\mathbb C)\)，\(t>0\)，并定义 \(Mf(x)=\sup_{r>0}|B(x,r)|^{-1}\int_{B(x,r)}|f|\)，其中球为开球、集合大小取 Lebesgue 测度。则
\[
 |\{Mf>t\}|\leq \dfrac{5^d}{t}\|f\|_1.
\]
\end{theorem}
\begin{proof}
记 \(E=\{Mf>t\}\)，并先固定紧集 \(K\subset E\)。对每个 \(x\in K\)，由上确界严格大于 \(t\)，能选一个以 \(x\) 为中心、平均大于 \(t\) 的球。开球覆盖紧集，所以可取有限子覆盖。

在这有限族中选半径最大的球，将所有与它相交的球删除；在余下的球中重复。得到的不交球记为 \(B_j=B(c_j,r_j)\)。若 \(B(a,s)\) 在选择 \(B_j\) 时被删，则 \(s\leq r_j\)、\(|a-c_j|<s+r_j\)。对其中任意点 \(y\)，有 \(|y-c_j|<s+(s+r_j)\leq3r_j\)。故被删球包含于 \(3B_j\)，所选球的三倍扩大已覆盖原来的有限族，也就覆盖 \(K\)。特别地，可以用 \(5B_j\) 作下列估计：
\[
 |K|\leq5^d\sum_j|B_j|\leq\dfrac{5^d}{t}\sum_j\int_{B_j}|f|\leq\dfrac{5^d}{t}\|f\|_1.
\]
第二步使用每个所选球的平均大于 \(t\)，最后一步才使用球互不相交。对全部紧 \(K\subset E\) 取上确界，Lebesgue 内正则性给出 \(|E|\) 的同一上界，得到结论。即使最初不知道 \(|E|\) 有限，这种内逼近仍然适用。证明只进行了有限族的选择，不需要预设能从无穷族中选出最大半径球。
\end{proof}
注意估计控制的是超水平集合的测度，没有直接断言 \(Mf\) 可积。例如 \(f=\mathbf1_{B(0,1)}\)。在 \(|x|\geq2\) 时，球 \(B(x,|x|+1)\) 包含单位球，故 \(Mf(x)\geq c_d|x|^{-d}\)。反过来，任何碰到单位球的中心球半径至少为 \(|x|-1\)，其所含质量又不超过 \(|B(0,1)|\)，所以 \(Mf(x)\leq C_d|x|^{-d}\)。尾部径向积分于是含有发散的 \(\int_2^\infty r^{-1}\,\mathrm dr\)。这个例子一面满足上述集合估计，一面否定一般的强 \(L^1\) 界，正是下一节引入弱空间的原因。

\subsection{中心/非中心极大函数}
\label{b2:subsec:170102}

对 \(f\in L^1_{\mathrm{loc}}(\mathbb R^d;\mathbb C)\)，若上确界改取所有包含 \(x\) 的开球，所得函数记为 \(M_u f(x)\)，称为\term[feizhongxin-jida-hanshu]{非中心极大函数}[uncentered maximal function]。这里“包含”不要求 \(x\) 是球心。给定 \(x\in B(a,r)\)，有 \(B(a,r)\subset B(x,2r)\)，故
\[\frac1{|B(a,r)|}\int_{B(a,r)}|f|
 \leq2^d\frac1{|B(x,2r)|}\int_{B(x,2r)}|f|.\]
对这类球取上确界，得到
\[
 Mf\leq M_u f\leq2^d Mf.
\]
两种平均的数值并不相同，但任何一者的范数或分布估计都可借此传给另一者，只改变维数常数。非中心函数也可测：\(M_uf(x)>t\) 当且仅当 \(x\) 位于某个平均大于 \(t\) 的开球内，所以该超水平集合是开球的并。对实直线上的区间输入，非中心区间可以把观察点放在端点附近，中心区间却要同时向两侧扩张；两者的差别来自允许怎样选择覆盖球。

\ProseParagraphBreak

对 \(R>0\)、\(f\in L^1_{\mathrm{loc}}(\mathbb R^d)\)，局部版本 \(M_Rf(x)\) 定义为 \(\sup_{0<r<R}|B(x,r)|^{-1}\int_{B(x,r)}|f|\)。若只考察紧集 \(K\)，选一个包含 \(K+B(0,R)\) 的有界球 \(U\)，令 \(F=f\cdot\mathbf1_U\)。则 \(F\in L^1\)，且 \(M_Rf(x)=M_RF(x)\leq MF(x)\) 对 \(x\in K\) 成立。这样才把全空间弱估计合法地用于局部函数。

\ProseParagraphBreak

局部可积本身不保证全尺度极大函数有限。例如在实直线上取 \(f(y)=|y|\)，则以原点为中心、半径 \(r\) 的平均为 \(r/2\)，故 \(Mf(0)=\infty\)。问题来自任意大的尺度；截断半径后，每个固定紧集只使用一个有界邻域内的积分。后面的局部微分论证也会按这个方式先截断位置。

\subsection{二进极大函数}
\label{b2:subsec:170103}

对 \(f\in L^1_{\mathrm{loc}}(\mathbb R^d;\mathbb C)\)，令 \(\mathcal D\) 为半开二进立方体 \(2^{-j}(k+[0,1)^d)\) 的族，其中 \(j\in\mathbb Z\)、\(k\in\mathbb Z^d\)。定义\term[erjin-jida-hanshu]{二进极大函数}[dyadic maximal function]
\[
 M_{\mathcal D}f(x)\coloneqq\sup_{Q\in\mathcal D,\,x\in Q}\dfrac1{|Q|}\int_Q|f|.
\]
\symidx{\(M_{\mathcal D}f\)：二进极大函数}
半开约定使每个尺度恰有一个立方体包含 \(x\)，同一格中的两个立方体不交或一者包含另一者。网格可数，故 \(M_{\mathcal D}f\) 可测。与球族相比，这里的尺度选择少了，却多了唯一的父立方体：每个 \(Q\) 都包含于一个边长加倍、体积为 \(2^d|Q|\) 的二进父块。下一章需要的恰是这个父子结构。

\Needspace{6\baselineskip}
\begin{proposition}[二进超水平集合分解]\label{b2:thm:17-dyadic-level}
设整数 \(d\geq1\)，\(f\in L^1(\mathbb R^d;\mathbb C)\)、\(t>0\)。固定标准半开二进网格 \(\mathcal D=\{\,2^{-j}(k+[0,1)^d):j\in\mathbb Z,\ k\in\mathbb Z^d\,\}\)，定义
\[
 M_{\mathcal D}f(x)=\sup_{Q\in\mathcal D,\,x\in Q}\frac1{|Q|}\int_Q|f|.
\]
则 \(\{M_{\mathcal D}f>t\}\) 是一族互不相交的极大二进立方体 \(Q_j\) 的并，且
\[
 t<\dfrac1{|Q_j|}\int_{Q_j}|f|\leq2^dt,
 \quad \sum_j|Q_j|\leq t^{-1}\|f\|_1.
\]
\end{proposition}
\begin{proof}
对平均大于 \(t\) 的任意二进块 \(Q\)，其第 \(n\) 个祖先 \(Q^{(n)}\) 的平均满足
\[\frac1{|Q^{(n)}|}\int_{Q^{(n)}}|f|
 \leq\frac{\|f\|_1}{2^{nd}|Q|}\longrightarrow0.\]
所以只有有限多个祖先仍可能平均大于 \(t\)；取这些超水平祖先中最大的一个。注意不能仅在第一次遇到平均不超过 \(t\) 的父块时停下，因为更大的祖先可能又包含别处的质量。所有足够大的祖先都不超水平，才保证极大块真正存在。

取全部极大超水平块。二进嵌套性说明两个不同的极大块必不交；网格可数，故它们组成有限或可数族 \(Q_j\)。\(M_{\mathcal D}f(x)>t\) 的每一点都在某个超水平块中，再包含于一个极大块，反向包含也由定义成立。每个 \(Q_j\) 的父块不超水平，因此
\[\int_{Q_j}|f|\leq\int_{\operatorname{parent}(Q_j)}|f|
 \leq t\,2^d|Q_j|.\]
这给出平均上界；下界来自超水平条件。最后由不交性，
\[t\sum_j|Q_j|\leq\sum_j\int_{Q_j}|f|\leq\|f\|_1,\]
得到测度估计。
\end{proof}
这项分解还记录坏区域的平均上界，是下一章分解函数的依据。仅有球覆盖弱型估计并不能直接提供这样的父子关系。

\ProseParagraphBreak

单个固定网格并不等同于全部球平均。对 \(x\in Q\)，以 \(x\) 为中心、半径 \(\sqrt d\,\ell(Q)\) 的球包含 \(Q\)，所以 \(M_{\mathcal D}f\leq C_dMf\)。反向的逐点比较对一个网格一般失败：在实直线上取 \(f=\mathbf1_{(0,1)}\)、\(x=-1\)，包含 \(-1\) 的标准二进区间都留在负半轴，因而 \(M_{\mathcal D}f(-1)=0\)；中心球却可跨过原点取得正平均。二进算子的弱型估计应由它自身的分解证明，而不能把两个算子直接当作数值等价的函数。

\subsection{极大函数微分法}
\label{b2:subsec:170104}

平均能恢复函数，并非因为任意可测代表在每一点都连续，而是坏点的测度可由逼近误差控制。
\Needspace{6\baselineskip}
\begin{proposition}[Lebesgue 微分的极大函数证明]\label{b2:thm:17-differentiation}
设整数 \(d\geq1\)，\(f\in L^1_{\mathrm{loc}}(\mathbb R^d;\mathbb C)\)，并选取一个有限值可测代表。则对几乎处处 \(x\in\mathbb R^d\)，有
\[
 \lim_{r\downarrow0}\dfrac1{|B(x,r)|}\int_{B(x,r)}|f(y)-f(x)|\,\mathrm dy=0.
\]
另令 \(Q_j(x)\) 为标准半开二进网格中唯一包含 \(x\)、边长 \(2^{-j}\) 的立方体，则还对几乎处处 \(x\) 有
\[\lim_{j\to\infty}\frac1{|Q_j(x)|}\int_{Q_j(x)}|f(y)-f(x)|\,\mathrm dy=0.\]
\end{proposition}
\begin{proof}
先设 \(f\in L^1\)。连续紧支集函数在 \(L^1(\mathbb R^d)\) 中稠密（第一卷命题~10.2.4），故对任意 \(\varepsilon>0\)，可取 \(g\in C_c(\mathbb R^d)\)，使 \(\|f-g\|_1<\varepsilon\)。令 \(Df(x)\) 为球平均绝对偏差的上极限。它是可测函数：固定半径的参数积分可测，半径的连续性又允许在有理半径上取上极限。由
\[
 |f(y)-f(x)|\leq |f(y)-g(y)|
       +|g(y)-g(x)|+|g(x)-f(x)|
\]
及 \(g\) 的连续性，得到
\[Df(x)\leq M(f-g)(x)+|f(x)-g(x)|.\]
因此对 \(t>0\)，由极大弱型和 Chebyshev 不等式，
\[
 |\{Df>2t\}|\leq(5^d+1)\varepsilon/t.
\]
左侧集合由 \(f,t\) 固定，不随所选 \(g\) 改变。令 \(\varepsilon\downarrow0\)，可知对每个固定 \(t>0\) 该集合为零测集，再只对 \(t=1/n\) 取可数并，得到 \(Df=0\) 几乎处处。

二进版本令 \(D_{\mathcal D}f\) 为所列平均偏差的上极限。因为 \(\operatorname{diam}Q_j(x)=\sqrt d\,2^{-j}\to0\)，连续函数的平均偏差仍趋零，所以
\[D_{\mathcal D}f\leq M_{\mathcal D}(f-g)+|f-g|.\]
二进超水平分解给出常数为一的弱估计，其余误差集合论证相同。

最后设 \(f\) 仅局部可积。对每个整数 \(N\geq1\)，置 \(f_N=f\cdot\mathbf1_{B(0,N+1)}\in L^1\)。在 \(B(0,N)\) 内，各点的足够小球及足够小二进块都落在 \(B(0,N+1)\) 内，其平均偏差与 \(f_N\) 的相同。于是可在 \(B(0,N)\) 上应用已证结论，再合并整数 \(N\) 对应的可数个零测例外集。
\end{proof}
因此在二进分解的好区域上可以把“所有小立方体平均不超过 \(t\)”转为“\(|f|\leq t\) 几乎处处”。这一步将在\secref{b2:sec:1802}使用。

% !TeX root = ../../Book2.tex
\section{弱 \texorpdfstring{\(L^p\)}{Lᵖ} 空间}
\label{b2:sec:1702}

上一节的极大函数可能不可积，却仍满足 \(|\{Mf>t\}|\leq C\|f\|_1/t\)。这不是没有得到任何控制，而是控制的对象变成了“超过某个高度的区域有多大”。本节先把这种信息写成分布函数，再说明它怎样决定强、弱 \(L^p\) 以及 Lorentz 空间。分布和重排的基本恒等式适用于一般测度空间；局部可积性、卷积和简单函数逼近的具体应用会明确回到 Lebesgue 空间。

\subsection{分布函数}
\label{b2:subsec:170201}

设 \((X,\mathcal A,\mu)\) 为测度空间，\(f:X\to\mathbb C\) 可测，其\term[fenbu-hanshu]{分布函数}[distribution function]定义为
\[d_f(t)\coloneqq\mu\{\,x\in X:|f(x)|>t\,\},\quad t\geq0.\]
\symidx{\(d_f(t)\)：函数的分布函数}
它允许取无穷，且只保留超水平集合的测度，不保留集合的位置。它单调不增；对 \(t_n\downarrow t\)，集合 \(\{|f|>t_n\}\) 递增到 \(\{|f|>t\}\)，所以 \(d_f(t_n)\to d_f(t)\)，即右连续。这里使用从下连续性，无须先假设某个超水平集合有限。

\ProseParagraphBreak

把 \(|f(x)|\) 想成竖直线段的长度，就有
\[|f(x)|=\int_0^\infty\mathbf1_{\{|f(x)|>t\}}\,\mathrm dt.\]
先累加各点的竖直长度，与先累加每一层的水平测度，得到同一个总量。对幂次作相同分层，产生下面的公式。

\Needspace{6\baselineskip}
\begin{proposition}[层蛋糕公式]\label{b2:thm:17-layer-cake}
设 \((X,\mathcal A,\mu)\) 为测度空间，\(h:X\to[0,\infty]\) 可测，\(d_h(t)=\mu\{h>t\}\)，\(0<p<\infty\)。则有允许两端同时为无穷的恒等式
\[\int_Xh^p\,\mathrm d\mu=p\int_0^\infty t^{p-1}d_h(t)\,\mathrm dt.\]
特别地，对实或复可测函数 \(f\) 取 \(h=|f|\)，便得到 \(\int_X|f|^p\,\mathrm d\mu=p\int_0^\infty t^{p-1}d_f(t)\,\mathrm dt\)。
\end{proposition}
\begin{proof}
先设 \(h\) 为非负简单函数，将其正值写在有限个不交可测集合 \(E_j\) 上，记这些值为 \(a_j>0\)。两侧都等于 \(\sum_j a_j^p\mu(E_j)\)：右侧只需对有限个阶梯函数分别积分。若某个正值集合测度无穷，则两侧都为无穷，等式仍成立。

一般情形取非负简单函数 \(h_n\uparrow h\)。对每个 \(t\geq0\)，集合 \(\{h_n>t\}\) 递增到 \(\{h>t\}\)，故 \(d_{h_n}(t)\uparrow d_h(t)\)。分别在 \(X\) 和高度轴上用单调收敛，把简单函数等式传到极限即可。这一证明不要求 \(X\) 为 \(\sigma\)-有限；若改用积空间 Tonelli，必须另外检查所采用版本的假设。
\end{proof}
例如 \(f=a\cdot\mathbf1_E\)、\(a>0\)、\(\mu(E)=m<\infty\) 时，\(d_f(t)=m\) 于 \(0\leq t<a\)，其余为零。公式给出 \(p\int_0^a t^{p-1}m\,\mathrm dt=a^pm\)，与直接积分一致。一般函数可以在高水平有尖峰，也可以在低水平有长尾；在无穷测度空间上，这两端都必须检查。

\subsection{弱型 \texorpdfstring{\(L^p\)}{Lp}}
\label{b2:subsec:170202}

Chebyshev 不等式说明，强 \(L^p\) 函数满足 \(d_f(t)\leq t^{-p}\|f\|_p^p\)。反向若只知道这个逐高度上界，层蛋糕中的被积函数却被估成 \(Ct^{-1}\)，一般无法积分。因而应把“每个高度的集合都不太大”与“所有高度的总积分有限”区分开来。

\ProseParagraphBreak

设 \((X,\mathcal A,\mu)\) 为测度空间，\(f:X\to\mathbb C\) 可测，\(0<p<\infty\)。若
\[\|f\|_{p,\infty}\coloneqq\sup_{t>0}t\cdot d_f(t)^{1/p}<\infty,\]
则称 \(f\) 属于\term[ruo-lp-kongjian]{弱 \(L^p\) 空间}[weak Lp space]，记为 \(L^{p,\infty}(X,\mu)\)。该量等于使 \(d_f(t)\leq(A/t)^p\) 对全部 \(t>0\) 成立的最小 \(A\)。它是拟范数：若 \(A=\|f\|_{p,\infty}\)、\(B=\|g\|_{p,\infty}\)，则超水平集合的包含关系给出
\[\|f+g\|_{p,\infty}\leq2(A^p+B^p)^{1/p}\leq C_p(A+B).\]
常数一的三角不等式确实可能失败。在两个互不相交、测度均为一的集合 \(E,F\) 上，取 \(f=2\cdot\mathbf1_E+\mathbf1_F\)、\(g=\mathbf1_E+2\cdot\mathbf1_F\)。二者的弱 \(L^1\) 拟范数都是 \(2\)，而和等于 \(3\cdot\mathbf1_{E\cup F}\)，弱拟范数为 \(6\)。

\Needspace{6\baselineskip}
\begin{proposition}[弱范数与有限集合积分]\label{b2:thm:17-weak-inclusions}
设 \((X,\mathcal A,\mu)\) 为测度空间，\(f:X\to\mathbb C\) 可测，\(0<p<\infty\)，并令 \(\|f\|_{p,\infty}=\sup_{t>0}t\cdot\mu\{|f|>t\}^{1/p}\)。则 \(\|f\|_{p,\infty}\leq\|f\|_{L^p(X,\mu)}\)。若 \(A=\|f\|_{p,\infty}<\infty\)、\(0<r<p\)，且 \(E\) 可测、\(0<\mu(E)<\infty\)，则
\[\|f\|_{L^r(E)}\leq
 \left(\frac p{p-r}\right)^{1/r}
 A\,\mu(E)^{1/r-1/p}.\]
特别地，\(p>1\) 时有
\[\int_E|f|\,\mathrm d\mu\leq p'\|f\|_{p,\infty}\cdot\mu(E)^{1/p'},
 \quad p'=\frac p{p-1}.\]
反之，设 \(1<p<\infty\)、\((X,\mathcal A,\mu)\) 半有限；若存在 \(C<\infty\)，使每个有限测度可测集 \(E\) 都满足 \(\int_E|f|\,\mathrm d\mu\leq C\mu(E)^{1/p'}\)，则 \(f\in L^{p,\infty}\) 且 \(\|f\|_{p,\infty}\leq C\)。
\end{proposition}
\begin{proof}
第一项由 Chebyshev 得到。记 \(m=\mu(E)>0\)。若 \(A=0\)，则 \(f=0\) 几乎处处，结论直接成立。否则限制在 \(E\) 上的分布函数不超过 \(\min\{m,A^pt^{-p}\}\)。这两项在 \(a=Am^{-1/p}\) 处相等，故应在这个高度分开积分。层蛋糕给出
\[
 \begin{aligned}
 \int_E|f|^r\,\mathrm d\mu
 &\leq r\int_0^a t^{r-1}m\,\mathrm dt
       +rA^p\int_a^\infty t^{r-p-1}\,\mathrm dt\\
 &=m a^r+\frac{rA^p}{p-r}a^{r-p}
  =\frac p{p-r}A^rm^{1-r/p}.
 \end{aligned}
\]
这里 \(r<p\) 正好保证第二项在无穷端收敛。取 \(r\) 次方根即得；\(r=1<p\) 给出所列集合积分界。零测度集合的积分另为零。

证明反向，固定 \(t>0\)，取 \(F\subset\{|f|>t\}\) 且 \(0<\mu(F)<\infty\)。假设给出 \(t\mu(F)\leq C\mu(F)^{1/p'}\)，从而 \(\mu(F)\leq(C/t)^p\)。半有限性保证任意可测集合的测度等于其中有限测度子集测度的上确界，故 \(d_f(t)\leq(C/t)^p\)。这也解释了为什么反向需要该测度条件；Lebesgue 测度满足它。
\end{proof}
在 Lebesgue 空间上，\(p>1\) 的弱 \(L^p\) 函数由此局部可积。\(p=1\) 时不能照搬：例如 \(|x|^{-d}\cdot\mathbf1_{\{0<|x|<1\}}\) 属于弱 \(L^1\)，却在原点附近不可积。另一方面，降低到任意 \(0<r<p\) 的结论仍有效；严格不等号来自上面看得见的高度积分。

\ProseParagraphBreak

设 \((X,\mu),(Y,\nu)\) 为测度空间，\(1\leq p,q\leq\infty\)，\(T\) 把 \(L^p(X,\mu)\) 输入映为 \(Y\) 上的可测函数。若 \(\|Tf\|_{q,\infty}\leq A\|f\|_p\)，称它具有\term[ruoxing-guji]{弱型估计}[weak type estimate] \((p,q)\)，强型则把目标换成 \(L^q\)。当 \(q=\infty\) 时直接使用本性上确界。例如上一节的中心极大算子具有弱型 \((1,1)\)，但单位球输入的远场计算否定了强型 \((1,1)\)。

\ProseParagraphBreak

对整数 \(d\geq1\)、\(1<p<\infty\)，在 \(\mathbb R^d\) 上定义
\[N_p(f)\coloneqq\sup_{0<|E|<\infty}|E|^{-1/p'}\int_E|f|,\]
其中 \(E\) 可测。上一命题的两个方向给出
\[\|f\|_{p,\infty}\leq N_p(f)\leq p'\|f\|_{p,\infty}.\]
\(N_p\) 的三角不等式直接来自集合积分的三角不等式，零范数也只对应零等价类。因此它提供一个等价的真正范数。上确界可以只取有界集合，因为 \(E\cap B(0,n)\uparrow E\)，积分和测度都按单调收敛逼近。这种集合积分写法还适合卷积：它不要求原输入在全空间可积。

\Needspace{6\baselineskip}
\begin{proposition}[弱空间的完备性与卷积]\label{b2:prop:17-weak-complete}
设整数 \(d\geq1\)、\(1<p<\infty\)，\(p'=p/(p-1)\)。实或复函数空间 \(L^{p,\infty}(\mathbb R^d)\) 关于范数
\[N_p(f)=\sup_{\substack{E\subset\mathbb R^d\ \text{可测}\\0<|E|<\infty}}
 |E|^{-1/p'}\int_E|f|\]
完备。若 \(k\in L^1(\mathbb R^d)\)、\(f\in L^{p,\infty}(\mathbb R^d)\) 均为实或复函数，则 \(k*f\) 由几乎处处绝对收敛的积分定义，且
\[N_p(k*f)\leq\|k\|_1\cdot N_p(f).\]
\end{proposition}
\begin{proof}
设 \(f_n\) 在 \(N_p\) 中 Cauchy。对每个有限球 \(B\)，有
\[\|f_n-f_m\|_{L^1(B)}\leq |B|^{1/p'}\cdot N_p(f_n-f_m).\]
\(L^1(B)\) 的完备性给出局部极限；不同球上的极限因几乎处处唯一而相容，故拼成局部可积函数 \(f\)。沿整数半径球作对角抽取，可得到一个在全空间几乎处处趋于 \(f\) 的子列。

给定 \(\varepsilon>0\)，先取 \(N\)，使 \(N_p(f_n-f_m)\leq\varepsilon\) 对 \(n,m\geq N\) 成立。固定 \(n\geq N\)，对任意有限测度集合 \(E\) 沿上述子列使用 Fatou，便得
\[\int_E|f_n-f|\leq\varepsilon|E|^{1/p'}.\]
此界对全部 \(E\) 使用同一个 \(N\)，所以 \(N_p(f_n-f)\leq\varepsilon\)。再由 \(N_p\) 的三角不等式知 \(N_p(f)<\infty\)，完备性成立。

对卷积，欧氏空间上的 Lebesgue 积测度空间是 \(\sigma\)-有限的。Tonelli 和平移不变性给出
\[
 \begin{aligned}
 \int_E\int |k(y)|\cdot|f(x-y)|\,\mathrm dy\,\mathrm dx
 &=\int |k(y)|\int_{E-y}|f(z)|\,\mathrm dz\,\mathrm dy\\
 &\leq |E|^{1/p'}\|k\|_1\cdot N_p(f).
 \end{aligned}
\]
取整数半径球 \(E\)，有限右端说明内积分在全空间几乎处处有限，故定义了绝对收敛的卷积。它控制 \(|k*f|\)，再对全部有限测度 \(E\) 取上确界即得所需范数界。
\end{proof}
完备性不意味着有界或光滑紧支集函数在弱拟范数中稠密。设 \(d\geq1\)、\(0<p<\infty\)，令 \(f(x)=|x|^{-d/p}\cdot\mathbf1_{\{0<|x|<1\}}\)，并记 \(\omega_d=|B(0,1)|\)。当 \(t\geq1\) 时，\(d_f(t)=\omega_dt^{-p}\)，故 \(f\in L^{p,\infty}\)。对任意有界 \(g\)，在 \(t>\max\{1,\|g\|_\infty\}\) 时，\(\{|f|>2t\}\subset\{|f-g|>t\}\)，所以
\[\|f-g\|_{p,\infty}\geq\frac12\omega_d^{1/p}.\]
无论有界截断的高度多大，这个误差都不趋零。处理弱空间上的算子延拓时，必须使用已证明的积分定义或另选收敛方式，不能直接套用强 \(L^p\) 的稠密性。

\subsection{Lorentz 空间初步}
\label{b2:subsec:170203}

弱拟范数只取 \(t\cdot d_f(t)^{1/p}\) 的最高值，没有累加不同高度的贡献。为了区分同样临界的幂次尾部中是否还带有可积的对数改善，可以保留整个分布，而不只取一个上确界。先把函数值按大小排列：设 \((X,\mathcal A,\mu)\) 为测度空间、\(f:X\to\mathbb C\) 可测，定义其\term[dijian-chongpai]{递减重排}[decreasing rearrangement]
\[f^*(s)\coloneqq\inf\{\,t\geq0:d_f(t)\leq s\,\},\quad s>0,\]
空集的下确界取无穷。变量 \(s\) 表示测度大小，不是原空间的位置；这里的 \(f^*\) 也不表示微分理论中的点态代表。它是分布函数的广义逆：寻找使超水平集合缩到不超过 \(s\) 的最低高度。

\Needspace{6\baselineskip}
\begin{proposition}[重排与等分布]\label{b2:prop:17-rearrangement}
设 \((X,\mathcal A,\mu)\) 为测度空间，\(f:X\to\mathbb C\) 可测，\(d_f(t)=\mu\{|f|>t\}\)、\(t\geq0\)，并令 \(f^*(s)=\inf\{t\geq0:d_f(t)\leq s\}\)、\(s>0\)。则 \(f^*\) 单调不增，且对每个 \(t\geq0\)，
\[|\{\,s>0:f^*(s)>t\,\}|=d_f(t),\]
左侧取半直线的 Lebesgue 测度。因而对 \(0<p<\infty\)，
\[\int_X|f|^p\,\mathrm d\mu=\int_0^\infty f^*(s)^p\,\mathrm ds,\quad
 \|f\|_{p,\infty}=\sup_{s>0}s^{1/p}f^*(s),\]
各式允许无穷值。
\end{proposition}
\begin{proof}
增大 \(s\) 会扩大取下确界的集合，故 \(f^*\) 单调不增。若 \(d_f(t)\leq s\)，则 \(t\) 已在该集合中，\(f^*(s)\leq t\)。反之，若 \(d_f(t)>s\)，分布函数的右连续性保证存在 \(\delta>0\)，使 \(d_f(t+\delta)>s\)，于是 \(f^*(s)\geq t+\delta>t\)。因此
\[f^*(s)>t\quad\Longleftrightarrow\quad s<d_f(t).\]
这给出等分布恒等式；分别对 \(|f|\) 和 \(f^*\) 使用层蛋糕即得积分等式。

对弱拟范数，固定 \(s\)，让 \(0<t<f^*(s)\) 趋向 \(f^*(s)\)，上述等价关系给出 \(s^{1/p}f^*(s)\leq\|f\|_{p,\infty}\)。反向固定 \(t>0\)，让 \(0<s<d_f(t)\) 趋向 \(d_f(t)\)，得到反向上界。无穷值时按任意大的有限 \(s,t\) 解释同一论证。
\end{proof}
即使原空间有原子，这个等分布结论也成立；重排是在另一个连续测度坐标上表示原分布，并未要求把原子拆开。它也说明，两函数若分布相同，其强、弱范数相同，但这并不表示两函数在原空间上接近。

\ProseParagraphBreak

例如在 \(\mathbb R^d\) 取两个不交可测集 \(E,F\)，\(|E|=|F|=1\)，令 \(f=2\cdot\mathbf1_E+\mathbf1_F\)。\cref{b2:tab:17-distribution-rearrangement}将先前的两层模型按高度和测度两种坐标写出。原来的形状可以很不规则，重排却只留下两段水平阶梯。
\begin{table}[htbp]
\centering
\caption[两层函数的分布与重排]{两层函数的分布与重排。\(E,F\subset\mathbb R^d\) 不交且测度均为一，\(f=2\cdot\mathbf1_E+\mathbf1_F\)；\(t\) 是高度，\(s\) 是重排后的测度坐标。}
\label{b2:tab:17-distribution-rearrangement}
\begin{tabularx}{\textwidth}{@{}*{4}{>{\centering\arraybackslash}X}@{}}
\toprule
高度区间 & \(d_f(t)\) & 测度坐标区间 & \(f^*(s)\)\\
\midrule
\(0\leq t<1\) & \(2\) & \(0<s<1\) & \(2\)\\
\(1\leq t<2\) & \(1\) & \(1\leq s<2\) & \(1\)\\
\(t\geq2\) & \(0\) & \(s\geq2\) & \(0\)\\
\bottomrule
\end{tabularx}
\end{table}

对 \(0<p<\infty\)、\(0<q\leq\infty\)，\term[lorentz-kongjian]{Lorentz 空间}[Lorentz space]由下列拟范数定义：
\[
 \|f\|_{L^{p,q}}\coloneqq
 \left(\int_0^\infty[s^{1/p}f^*(s)]^q\frac{\mathrm ds}{s}\right)^{1/q}
 \quad(0<q<\infty),
\]
\(q=\infty\) 时取 \(\sup_{s>0}s^{1/p}f^*(s)\)。主指数 \(p\) 固定幂次尺度，第二指数 \(q\) 记录各测度尺度的累加方式。\(\mathrm ds/s\) 对伸缩 \(s\mapsto as\) 不变，因此每一倍增尺度受到同样对待。由重排命题，\(L^{p,p}=L^p\) 且范数相等，\(L^{p,\infty}\) 也与前面定义完全一致。
\symidx{\(f^*(s)\)：递减重排}
对表中的两层函数，可直接算得
\[
 \|f\|_{L^{p,q}}^q=\frac pq(2^q+2^{q/p}-1)\quad(q<\infty),\quad
 \|f\|_{p,\infty}=\max\{2,2^{1/p}\}.
\]
令 \(q=p\) 便恢复 \(\|f\|_p^p=2^p+1\)。这些等式展示的是同一个输入的不同大小信息，而非几种互不相干的空间定义。

\ProseParagraphBreak

固定 \(0<p<\infty\)，较小的第二指数确实给出较强的条件。对 \(0<q_0<\infty\)，单调性给出
\[
 \|f\|_{L^{p,q_0}}^{q_0}
 \geq f^*(s)^{q_0}\int_0^s r^{q_0/p-1}\,\mathrm dr
 =\frac p{q_0}[s^{1/p}f^*(s)]^{q_0}.
\]
故归一化重排 \(H(s)=s^{1/p}f^*(s)\) 的上确界受该拟范数控制。若 \(q_0<q_1<\infty\)，再用
\[\int_0^\infty H^{q_1}\frac{\mathrm ds}{s}
 \leq\|H\|_\infty^{q_1-q_0}\int_0^\infty H^{q_0}\frac{\mathrm ds}{s},\]
便得 \(L^{p,q_0}\subset L^{p,q_1}\)；\(q_1=\infty\) 已由上确界界涵盖。不能仅把它视为无穷测度空间上普通 \(L^q\) 的包含关系，关键额外信息是 \(f^*\) 的单调性。

\ProseParagraphBreak

临界幂次和对数使第二指数的作用更清楚。先设 \(a>0\)、\(0<p<\infty\)。径向积分给出
\[
 \begin{aligned}
 |x|^{-a}\cdot\mathbf1_{\{0<|x|<1\}}\in L^p
 &\quad\Longleftrightarrow\quad ap<d,\\
 |x|^{-a}\cdot\mathbf1_{\{|x|>1\}}\in L^p
 &\quad\Longleftrightarrow\quad ap>d.
 \end{aligned}
\]
在 \(ap=d\) 时，两者都只有临界弱界；相加所得纯幂次 \(|x|^{-a}\) 属于 \(L^{d/a,\infty}\)，却不属于任何全空间强 \(L^p\)。前一种失败来自高水平的原点奇性，后一种来自低水平的远尾。

\ProseParagraphBreak

再固定 \(0<p<\infty\)、\(\beta\in\mathbb R\)，在足够小的 \(0<s<s_0\) 上取
\[F(s)=s^{-1/p}(\log(e/s))^{-\beta},\]
并在 \(s\geq s_0\) 置零。选择 \(s_0\) 足够小使此函数在 \((0,s_0)\) 单调不增；例如 \(\log(e/s_0)>p\max\{\beta,0\}+1\) 即可。因此这个半直线上的模型本身就是自己的重排。对 \(q<\infty\)，代入定义并置 \(v=\log(e/s)\)，得到
\[\|F\|_{L^{p,q}}^q=\int_{\log(e/s_0)}^\infty v^{-\beta q}\,\mathrm dv.\]
所以 \(F\in L^{p,q}\) 当且仅当 \(\beta q>1\)，而 \(F\in L^{p,\infty}\) 当且仅当 \(\beta\geq0\)。特别地，强 \(L^p\) 要求 \(\beta p>1\)。纯临界幂次对应 \(\beta=0\)，只留下上确界控制；加入足够的对数衰减，才让尺度积分收敛。

\ProseParagraphBreak

这些更细的输入条件能够与弱端点输出配对，恢复原本可能发散的积分。
\Needspace{6\baselineskip}
\begin{proposition}[Lorentz 端点配对]\label{b2:prop:17-lorentz-pairing}
设 \((X,\mathcal A,\mu)\) 为测度空间，\(1<p<\infty\)、\(p'=p/(p-1)\)，实或复可测函数 \(f\in L^{p,1}(X,\mu)\)、\(g\in L^{p',\infty}(X,\mu)\)。Lorentz 拟范数采用递减重排及测度 \(\mathrm ds/s\) 的上述定义。则乘积绝对可积，且
\[\int_X|fg|\,\mathrm d\mu
 \leq\|f\|_{L^{p,1}}\cdot\|g\|_{p',\infty}.\]
\end{proposition}
\begin{proof}
对可测集 \(E\)，将层蛋糕用于 \(|g|\cdot\mathbf1_E\)，并用等分布恒等式，得
\[
 \int_E|g|\,\mathrm d\mu
 \leq\int_0^\infty\min\{\mu(E),d_g(u)\}\,\mathrm du
 =\int_0^{\mu(E)}g^*(s)\,\mathrm ds.
\]
最后的等式是在两个 Lebesgue 半直线之间使用非负积分换序，允许 \(\mu(E)=\infty\)。接着对测度 \(\nu(E)=\int_E|g|\,\mathrm d\mu\) 应用层蛋糕于 \(|f|\)，得到
\[
 \begin{aligned}
 \int_X|fg|\,\mathrm d\mu
 &=\int_0^\infty\int_{\{|f|>t\}}|g|\,\mathrm d\mu\,\mathrm dt\\
 &\leq\int_0^\infty\int_0^{d_f(t)}g^*(s)\,\mathrm ds\,\mathrm dt
 =\int_0^\infty f^*(s)g^*(s)\,\mathrm ds.
 \end{aligned}
\]
最后一步再用 \(s<d_f(t)\) 当且仅当 \(t<f^*(s)\)，只交换两个非负 Lebesgue 积分。代入 \(g^*(s)\leq\|g\|_{p',\infty}s^{-1/p'}\)，右端即为不超过 \(\|g\|_{p',\infty}\cdot\|f\|_{L^{p,1}}\) 的有限数。
\end{proof}
分层时需要注意恢复的对象：一般实、复函数的层蛋糕恢复的是 \(|f|\)，并非带符号或相位的 \(f\)。上述证明先控制绝对乘积，所以没有丢失相位条件。

\ProseParagraphBreak

这里一般不能把 \(L^{p,1}\) 放宽为 \(L^p\)。在 \((0,s_0)\) 上取同位置的两个函数
\[
 f(s)=s^{-1/p}(\log(e/s))^{-\beta},\quad g(s)=s^{-1/p'},
 \quad 1/p<\beta\leq1.
\]
前者是上面的对数模型，属于 \(L^p\)；后者属于弱 \(L^{p'}\)，但
\[\int_0^{s_0}|fg|\,\mathrm ds
 =\int_0^{s_0}s^{-1}(\log(e/s))^{-\beta}\,\mathrm ds=\infty.\]
若改为 \(\beta>1\)，则 \(f\in L^{p,1}\)，配对命题保证积分收敛。这一差别正好来自高度尺度上多一次可积求和。

\ProseParagraphBreak

最后考虑一种常见的输入端点：只先对集合示性函数检验算子。这样的\term[xianzhi-ruoxing-guji]{限制弱型估计}[restricted weak type estimate]比全部输入的弱型估计少一层要求，却自然产生 Lorentz 输入。
\Needspace{6\baselineskip}
\begin{proposition}[限制弱型与 Lorentz 输入]\label{b2:prop:17-restricted-weak}
设整数 \(d,n\geq1\)，\(1\leq p<\infty\)、\(1<q<\infty\)、\(q'=q/(q-1)\)。设复线性算子 \(T\) 定义在 \(\mathbb R^d\) 上有限测度支集的简单函数等价类中，输出 \(\mathbb R^n\) 上可测函数等价类。若存在 \(A<\infty\)，使每个有限测度可测集 \(E\subset\mathbb R^d\) 满足
\[\|T\mathbf1_E\|_{q,\infty}\leq A|E|^{1/p},\]
则 \(T\) 有唯一连续线性延拓 \(L^{p,1}(\mathbb R^d)\to L^{q,\infty}(\mathbb R^n)\)，且
\[N_q(Tf)\leq C_{p,q}A\|f\|_{L^{p,1}},\quad
 N_q(h)=\sup_{0<|B|<\infty}|B|^{-1/q'}\int_B|h|.\]
所有集合和积分取相应维数的 Lebesgue 测度。目标连续性可等价地用 \(N_q\) 范数或弱 \(L^q\) 拟范数表示。
\end{proposition}
\begin{proof}
若 \(f\) 是非负简单函数，\(f=\int_0^\infty\mathbf1_{\{f>t\}}\,\mathrm dt\) 实际只是有限个示性函数的正线性组合。\(N_q\) 的三角不等式与 \(N_q\leq q'\|\cdot\|_{q,\infty}\) 给出
\[N_q(Tf)\leq q'A\int_0^\infty d_f(t)^{1/p}\,\mathrm dt
 =\frac{q'A}{p}\|f\|_{L^{p,1}}.\]
最后的等式来自等分布关系和两个非负 Lebesgue 积分的换序。复简单函数分成实部、虚部的正负四部分，每部分的 Lorentz 量不超过原量，所以常数至多增加四倍。

为使用连续延拓，还须证明这类简单函数在输入空间稠密。先记一个拟三角界：超水平集合包含关系给出
\[(f+g)^*(s)\leq f^*(s/2)+g^*(s/2),\]
因而换元 \(r=s/2\) 得
\[\|f+g\|_{L^{p,1}}\leq2^{1/p}
 (\|f\|_{L^{p,1}}+\|g\|_{L^{p,1}}).\]
这里的重排界可以先取略高于两个广义逆的高度，使两个超水平集合的测度各不超过 \(s/2\)，再令高度下降取得。

给定 \(f\in L^{p,1}(\mathbb R^d)\)，保留 \(|x|\leq N\)、\(1/N\leq|f|\leq N\) 的部分，记为 \(f_N\)，令 \(e_N=|f-f_N|\)。则 \(e_N\downarrow0\) 几乎处处，且 \(e_N^*\leq f^*\)。由前面第二指数的上确界控制，\(d_f(t)<\infty\) 对每个 \(t>0\) 成立。因此 \(d_{e_N}(t)\downarrow0\)，进而 \(e_N^*(s)\to0\) 对每个 \(s>0\) 成立。以可积函数 \(s^{1/p-1}f^*(s)\) 支配，得到 \(\|f-f_N\|_{L^{p,1}}\to0\)。每个 \(f_N\) 有界且支集测度有限，可在该集合上作一致的简单函数逼近；若误差至多 \(\varepsilon\cdot\mathbf1_E\)，其 Lorentz 量至多 \(p\varepsilon|E|^{1/p}\)。结合拟三角界，即得所需稠密性。

最后对简单函数逼近列应用已证估计于两项之差，输出在完备的 \(N_q\) 范数中 Cauchy，因而有极限。差估计保证极限不依赖逼近列，并保留线性和同一上界；稠密性又保证延拓唯一。
\end{proof}
这个命题解释了为什么仅检验集合时，输入自然需要 \(L^{p,1}\) 的高度累加，而不是只用 \(L^p\) 的单一幂积分。目标条件 \(q>1\) 用于前面建立的等价 Banach 范数及完备性；在 \(q=1\) 时，这个证明不能原样使用。下一节的插值则从全部输入的两个弱端点出发，利用随输出高度改变的大小分解。

% !TeX root = ../../Book2.tex
\section{Marcinkiewicz 插值}
\label{b2:sec:1703}

弱型估计只控制输出超过某一高度的集合。若直接把它代入层蛋糕公式，端点处常会出现发散的高度积分。可以先把输入按大小分成两部分：大值部分交给较低的可积指数，小值部分交给较高的可积指数。分割高度随正在考察的输出高度改变，两个端点就分别控制同一个高度积分的两端。

\subsection{次线性算子}
\label{b2:subsec:170301}

设 \((X,\mathcal A,\mu)\)、\((Y,\mathcal B,\nu)\) 为测度空间，\(\mathbb K=\mathbb R\) 或 \(\mathbb C\)，\(\mathcal D\) 是 \(X\) 上有限值可测 \(\mathbb K\)-值函数的几乎处处等价类组成的向量空间。设 \(T\) 把 \(\mathcal D\) 映为 \(Y\) 上有限值可测函数的几乎处处等价类。若对所有 \(f,g\in\mathcal D\)、\(c\in\mathbb K\)，有
\[
 |T(f+g)|\leq |Tf|+|Tg|,
 \quad |T(cf)|=|c|\cdot|Tf|
 \quad\text{几乎处处},
\]
则称 \(T\) 为\term[cixianxing-suanzi]{次线性算子}[sublinear operator]。极大算子在 \(L^1+L^\infty\) 等已保证输出几乎处处有限的定义域上满足这些条件；线性算子以及线性算子的模也满足这些条件。这里保留的是输出大小的三角不等式，不要求输出本身能够线性叠加。

\ProseParagraphBreak

为什么把大值交给较小的指数？设 \(1\leq p_0<p<p_1<\infty\)，\(f\in L^p(X,\mu)\)，并取 \(a>0\)。令
\[
 f^{>a}=f\cdot\mathbf1_{\{|f|>a\}},
 \quad f^{\leq a}=f\cdot\mathbf1_{\{|f|\leq a\}}.
\]
在大值集合上，\(p_0-p<0\) 给出 \(|f|^{p_0-p}\leq a^{p_0-p}\)；小值集合上则有 \(|f|^{p_1-p}\leq a^{p_1-p}\)。所以
\[
 \int_{|f|>a}|f|^{p_0}\,\mathrm d\mu
 \leq a^{p_0-p}\|f\|_p^p,
 \quad
 \int_{|f|\leq a}|f|^{p_1}\,\mathrm d\mu
 \leq a^{p_1-p}\|f\|_p^p.
\]
若 \(p_1=\infty\)，后一条换成 \(\|f^{\leq a}\|_\infty\leq a\)。这也证明了每个中间 \(L^p\) 输入都属于和空间 \(L^{p_0}+L^{p_1}\)：两个截断确实落在算子的端点定义域中。

\ProseParagraphBreak

以上粗界适合检查截断的可积性，暂时不要把它们代入整个插值证明。这样代入后，两项都只剩下 \(t^{-p}\|f\|_p^p\) 一类的分布上界，层蛋糕积分又会遇到 \(\int\mathrm dt/t\)。要保留比弱型更多的信息，必须让集合 \(\{|f|>a\}\) 和 \(\{|f|\leq a\}\) 留在高度积分里面。

\subsection{弱型估计}
\label{b2:subsec:170302}

回顾\secref{b2:sec:1702}的记号，\(d_h(t)=\nu(\{|h|>t\})\)、\(t>0\)，表示目标函数的分布函数。若 \(T\) 在 \(L^{p_i}(X)\) 上满足弱型 \((p_i,p_i)\)，其含义是存在 \(A_i\geq0\)，使
\[
 d_{Th}(t)\leq
 \left(\dfrac{A_i\|h\|_{L^{p_i}(X)}}{t}\right)^{p_i}
 \quad(h\in L^{p_i}(X),\ t>0).
\]
当 \(p_i=\infty\) 时，使用的是强型条件
\(\|Th\|_{L^\infty(Y)}\leq A_i\|h\|_{L^\infty(X)}\)。

\ProseParagraphBreak

对定义域中的分解 \(f=f_0+f_1\)，如果 \(|Tf_0|\leq t/2\) 且 \(|Tf_1|\leq t/2\)，次线性便给出 \(|Tf|\leq t\)。因此，除一个零测集外，
\[
 \{|Tf|>t\}\subseteq
 \{|Tf_0|>t/2\}\cup\{|Tf_1|>t/2\}.
\]
由此可以分别使用两个端点，而不需要 \(Tf=Tf_0+Tf_1\)。若较高端点是 \(L^\infty\)，且其常数 \(A_1>0\)，取 \(a=t/(2A_1)\) 就使小值部分的输出不超过 \(t/2\)；此时只须估计大值部分的超水平集合。

\subsection{Marcinkiewicz 插值}
\label{b2:subsec:170303}

\term[marcinkiewicz-chazhi]{Marcinkiewicz 插值定理}[Marcinkiewicz interpolation theorem]先把两个同指数的弱型端点连接起来。这是后面从弱 \((1,1)\) 与强 \((2,2)\) 推出奇异积分中间指数估计时使用的版本。

\Needspace{6\baselineskip}
\begin{theorem}[Marcinkiewicz]\label{b2:thm:17-marcinkiewicz}
设 \((X,\mathcal A,\mu)\)、\((Y,\mathcal B,\nu)\) 为测度空间，\(\mathbb K=\mathbb R\) 或 \(\mathbb C\)，\(1\leq p_0<p_1\leq\infty\)。设 \(T\) 定义在 \(L^{p_0}(X;\mathbb K)+L^{p_1}(X;\mathbb K)\) 上，输出 \(Y\) 上可测函数的几乎处处等价类，且满足 \(|T(f+g)|\leq|Tf|+|Tg|\)、\(|T(cf)|=|c|\cdot|Tf|\) 几乎处处。设 \(A_0,A_1\geq0\)，并假设对每个有限端点 \(p_i\)，有
\[
 \nu(\{|Th|>t\})\leq
 \left(\dfrac{A_i\|h\|_{L^{p_i}(X)}}{t}\right)^{p_i}
 \quad(h\in L^{p_i}(X),\ t>0);
\]
若 \(p_1=\infty\)，该端点改为
\(\|Th\|_{L^\infty(Y)}\leq A_1\|h\|_{L^\infty(X)}\)。
对 \(p_0<p<p_1\)，以 \(1/\infty=0\) 约定唯一的 \(0<\theta<1\)，使
\[
 \frac1p=\frac{1-\theta}{p_0}+\frac\theta{p_1}.
\]
则对所有 \(f\in L^p(X;\mathbb K)\)，有
\[
 \|Tf\|_{L^p(Y)}
 \leq C_{p,p_0,p_1}A_0^{1-\theta}A_1^\theta\|f\|_{L^p(X)},
\]
其中常数只依赖所列指数。源和目标空间不要求无原子或 \(\sigma\)-有限。
\end{theorem}
\begin{proof}
先设 \(A_0,A_1>0\)，固定 \(f\in L^p(X)\)。这里 \(p<\infty\)，所以
\[
 E=\{|f|>0\}=\bigcup_{n\geq1}\{|f|>1/n\},
 \quad \mu(\{|f|>1/n\})\leq n^p\|f\|_p^p<\infty.
\]
因此 \(\mu\) 在 \(E\) 上的限制是 \(\sigma\)-有限的。以下与输入有关的非负重积分都支撑在 \(E\times(0,\infty)\) 上，Tonelli 定理可在那里使用；目标上的层蛋糕公式则用\cref{b2:thm:17-layer-cake}的一般测度版本。

先讨论 \(p_1<\infty\)。取一个待选择的 \(\kappa>0\)，在输出高度 \(t\) 处按输入高度 \(a=\kappa t\) 分割。前面已证明两个截断分别属于 \(L^{p_0}\)、\(L^{p_1}\)，故超水平集合估计给出
\begin{equation}\label{b2:eq:17-height-split}
 \begin{aligned}
 d_{Tf}(t)\leq{}&\left(\frac{2A_0}{t}\right)^{p_0}
       \int_{|f|>\kappa t}|f|^{p_0}\,\mathrm d\mu\\
 &+\left(\frac{2A_1}{t}\right)^{p_1}
       \int_{|f|\leq\kappa t}|f|^{p_1}\,\mathrm d\mu.
 \end{aligned}
\end{equation}
选择 \(a\) 与 \(t\) 成正比，使固定输入值 \(s=|f(x)|\) 对应的高度区间恰为 \(0<t<s/\kappa\) 或 \(t\geq s/\kappa\)。两个端点因而分别控制高度零点和无穷远处。

将\eqref{b2:eq:17-height-split}乘以 \(pt^{p-1}\)，对 \(t>0\) 积分。层蛋糕公式与 Tonelli 给出
\[
 \begin{aligned}
 \|Tf\|_p^p\leq{}&p(2A_0)^{p_0}
   \int_E |f|^{p_0}
      \left(\int_0^{|f|/\kappa}t^{p-p_0-1}\,\mathrm dt\right)\mathrm d\mu\\
 &+p(2A_1)^{p_1}
   \int_E |f|^{p_1}
      \left(\int_{|f|/\kappa}^\infty t^{p-p_1-1}\,\mathrm dt\right)\mathrm d\mu\\
 ={}&p\left(
   \frac{(2A_0)^{p_0}\kappa^{p_0-p}}{p-p_0}
  +\frac{(2A_1)^{p_1}\kappa^{p_1-p}}{p_1-p}\right)\|f\|_p^p.
 \end{aligned}
\]
第一段内积分有限恰好用到 \(p>p_0\)，第二段有限用到 \(p<p_1\)。左端在计算开始时允许为无穷，右端有限才证明 \(Tf\in L^p\)。若先取 \(\kappa=1\)，已经得到所求强型估计。

为了同时看清端点常数的作用，令两个分子中的常数因子相等，取
\[
 \kappa=\left(\frac{(2A_0)^{p_0}}{(2A_1)^{p_1}}\right)^{1/(p_1-p_0)}.
\]
由倒数指数关系，
\[
 p(1-\theta)=\frac{p_0(p_1-p)}{p_1-p_0},
 \quad p\theta=\frac{p_1(p-p_0)}{p_1-p_0}.
\]
于是两个分子均为
\(2^pA_0^{p(1-\theta)}A_1^{p\theta}\)，并得到显式的上界
\[
 \|Tf\|_p\leq
 2\left[p\left(\frac1{p-p_0}+\frac1{p_1-p}\right)\right]^{1/p}
 A_0^{1-\theta}A_1^\theta\|f\|_p.
\]
这项选择使两个分子相等，从而显示两个端点常数在上界中的幂次。

若 \(p_1=\infty\)，取 \(a=t/(2A_1)\)。小值部分的输出不超过 \(t/2\)，因此
\[
 d_{Tf}(t)\leq\left(\frac{2A_0}{t}\right)^{p_0}
       \int_{|f|>t/(2A_1)}|f|^{p_0}\,\mathrm d\mu.
\]
换序后，固定 \(x\) 的高度积分为
\(\int_0^{2A_1|f(x)|}t^{p-p_0-1}\,\mathrm dt\)。故
\[
 \|Tf\|_p^p\leq
 \frac{p(2A_0)^{p_0}(2A_1)^{p-p_0}}{p-p_0}\|f\|_p^p,
\]
即
\[
 \|Tf\|_p\leq
 2\left(\frac p{p-p_0}\right)^{1/p}
 A_0^{p_0/p}A_1^{1-p_0/p}\|f\|_p.
\]
此时 \(\theta=1-p_0/p\)，仍是同一个几何幂。最后，若某个 \(A_i=0\)，先用正的较大常数 \(A_i+\varepsilon\) 作上述估计，再令 \(\varepsilon\downarrow0\)，便得到相应结论。
\end{proof}

两个高度积分解释了为什么结论只在开区间内成立：当 \(p=p_0\) 或有限的 \(p=p_1\) 时，其中一个内积分成为对数发散。改变分割常数可以调整算子范数的上界，不能消除这一端点障碍。

\ProseParagraphBreak

定理假设算子已经在和空间上定义，因而截断可以直接送入算子。另一种常见情形是线性算子先只定义在有限测度支集的简单函数上，并在那里满足同样的两个端点界。这个类对高度截断封闭，证明先给出该类的中间强型界；它在有限指数 \(L^p\) 中稠密，且
\(\|Tf_n-Tf_m\|_p=\|T(f_n-f_m)\|_p\leq C\|f_n-f_m\|_p\)，
所以目标 \(L^p\) 的完备性产生唯一连续延拓。对一般次线性算子，若另外具有
\(|Tf-Tg|\leq|T(f-g)|\)，也可作同一构造。非负值次线性算子具有这项差值控制：分别把 \(f=g+(f-g)\)、\(g=f+(g-f)\) 代入次可加性即可。若只知道本节定义中的模不等式，输出的相位变化并不受控，不能由此断言算子本身连续或有唯一延拓。

\ProseParagraphBreak

输入和输出的指数也可以不同。下面保留一个适用于局部势核的版本；除了层蛋糕公式，还需要用到指数至少为一时的积分三角不等式。

\Needspace{6\baselineskip}
\begin{proposition}\label{b2:prop:17-off-diagonal}
设 \((X,\mathcal A,\mu)\)、\((Y,\mathcal B,\nu)\) 为测度空间，\(\mathbb K=\mathbb R\) 或 \(\mathbb C\)，\(1\leq r<\infty\)。设 \(T\) 定义在 \(L^1(X;\mathbb K)+L^\infty(X;\mathbb K)\) 上，输出 \(Y\) 上可测函数的几乎处处等价类，满足 \(|T(f+g)|\leq|Tf|+|Tg|\)、\(|T(cf)|=|c|\cdot|Tf|\) 几乎处处。假设存在 \(A\geq0\)、\(B>0\)，使
\[
 \begin{aligned}
 \nu(\{|Th|>t\})&\leq
 \left(\frac{A\|h\|_{L^1(X)}}t\right)^r
 &&(h\in L^1(X),\ t>0),\\
 \|Th\|_{L^\infty(Y)}&\leq B\|h\|_{L^\infty(X)}
 &&(h\in L^\infty(X)).
 \end{aligned}
\]
则对 \(1<p<\infty\)、\(f\in L^p(X;\mathbb K)\)，有
\[
 \|Tf\|_{L^{pr}(Y)}
 \leq 2\left(\frac p{p-1}\right)^{1/(pr)}
 A^{1/p}B^{1-1/p}\|f\|_{L^p(X)}.
\]
这里不要求两个测度空间 \(\sigma\)-有限。
\end{proposition}
\begin{proof}
固定 \(f\in L^p(X)\)，其非零支集 \(E\) 仍为 \(\sigma\)-有限。对输出高度 \(t>0\)，取
\[
 f_0=f\cdot\mathbf1_{\{|f|>t/(2B)\}},
 \quad f_1=f-f_0.
\]
有 \(\|f_0\|_1\leq(t/(2B))^{1-p}\|f\|_p^p\)、\(\|f_1\|_\infty\leq t/(2B)\)，所以两个端点都能使用，且小值输出不超过 \(t/2\)。于是
\[
 d_{Tf}(t)\leq\left(\frac{2A}t\right)^r
 \left(\int_E|f(x)|\cdot\mathbf1_{\{t<2B|f(x)|\}}\,\mathrm d\mu(x)\right)^r.
\]
置 \(q=pr\)，按层蛋糕公式，
\[
 \|Tf\|_q^q\leq q(2A)^r
 \int_0^\infty t^{q-r-1}
 \left(\int_E |f(x)|\cdot\mathbf1_{\{t<2B|f(x)|\}}\,\mathrm d\mu(x)\right)^r\mathrm dt.
\]
与同指数证明相比，输入积分现在还带有 \(r\) 次幂，不能直接把这次幂移进 \(x\) 积分。

记 \(\mathrm d\rho(t)=t^{q-r-1}\,\mathrm dt\)，并将括号内的非负函数记为 \(H(t)\)。在这个高度测度上使用积分 Minkowski 不等式，得到
\[
 \begin{aligned}
 \|H\|_{L^r(\rho)}
 &\leq\int_E |f(x)|
       \left(\int_0^{2B|f(x)|}t^{q-r-1}\,\mathrm dt\right)^{1/r}\mathrm d\mu(x)\\
 &=\frac{(2B)^{p-1}}{(q-r)^{1/r}}\|f\|_p^p.
 \end{aligned}
\]
这里 \(q-r=r(p-1)>0\) 保证零点处可积。所用的积分不等式也可以直接核验：\(r=1\) 时它就是 Tonelli；\(r>1\) 时，令 \(r'=r/(r-1)\)，对任意非负 \(h\in L^{r'}(\rho)\)、\(\|h\|_{r'}\leq1\)，先由 Tonelli 交换 \(x,t\) 积分，再对每个 \(x\) 用 Hölder，便有
\[
 \int_0^\infty H(t)\cdot h(t)\,\mathrm d\rho(t)
 \leq\int_E |f(x)|
       \left(\int_0^{2B|f(x)|}\mathrm d\rho(t)\right)^{1/r}\mathrm d\mu(x).
\]
为避免预设 \(H\in L^r(\rho)\)，先取
\(H_n=\min\{H,n\}\cdot\mathbf1_{[1/n,n]}\)。它有有限 \(L^r(\rho)\) 范数；非零时用
\(h=H_n^{r-1}/\|H_n\|_r^{r-1}\) 测试，得到同一个范数上界，再由单调收敛令 \(n\to\infty\)。所有重积分都在 \(E\times(0,\infty)\) 这个 \(\sigma\)-有限乘积空间上进行。

将 \(H\) 的估计代回，取 \(1/r\) 次方，得到
\[
 \|Tf\|_q^p
 \leq q^{1/r}(2A)\cdot\frac{(2B)^{p-1}}{(q-r)^{1/r}}\|f\|_p^p
 =2^p\left(\frac p{p-1}\right)^{1/r}AB^{p-1}\|f\|_p^p.
\]
再取 \(1/p\) 次方即得结论。若 \(A=0\)，可先用正的较大弱型常数，再取极限。
\end{proof}

这个指数可以在参数线上读出：对端点 \((1,r)\) 与 \((\infty,\infty)\)，取 \(\theta=1-1/p\)，则
\[
 \frac1q=\frac{1-\theta}r=\frac1{pr},
 \quad A^{1-\theta}B^\theta=A^{1/p}B^{1-1/p}.
\]
高度积分中的 Minkowski 论证使这个参数计算成为强型估计；假如 \(r<1\)，该积分三角不等式不再成立，以上证明便不能照用。

\ProseParagraphBreak

例如在 \(\mathbb R^d\) 上取 \(0<\beta<d\)，令
\[
 k(x)=|x|^{-\beta}\cdot\mathbf1_{\{0<|x|<1\}},
 \quad r=\frac d\beta>1,
 \quad \omega_d=|B(0,1)|.
\]
球坐标积分与超水平集合的半径给出
\[
 \|k\|_1=\frac{d\omega_d}{d-\beta},
 \quad
 d_k(s)=\begin{cases}\omega_d,&0<s<1,\\
                         \omega_ds^{-r},&s\geq1,
          \end{cases}
 \quad \|k\|_{r,\infty}=\omega_d^{1/r}.
\]
由\secref{b2:sec:1702}的集合积分范数
\(N_r(h)=\sup_{0<|F|<\infty}|F|^{-1/r'}\int_F|h|\)，有
\(N_r(k)\leq r'\omega_d^{1/r}\)。对 \(h\in L^1\)，把\cref{b2:prop:17-weak-complete}用于可积因子 \(h\) 与弱空间因子 \(k\)，得
\[
 \|k*h\|_{r,\infty}\leq N_r(k*h)
 \leq N_r(k)\cdot\|h\|_1.
\]
另一方面，\(k\in L^1\) 给出
\(\|k*h\|_\infty\leq\|k\|_1\cdot\|h\|_\infty\)。卷积积分在 \(L^1+L^\infty\) 上具有兼容的定义：可积部分的卷积几乎处处绝对收敛，有界部分的卷积处处绝对收敛；其和等于原输入的卷积积分几乎处处。于是命题给出
\[
 \|k*f\|_{pr}
 \leq 2\left(\frac p{p-1}\right)^{1/(pr)}
       (r'\omega_d^{1/r})^{1/p}
       \left(\frac{d\omega_d}{d-\beta}\right)^{1-1/p}\|f\|_p.
\]
取 \(d=3\)、\(\beta=1\)、\(p=2\)，就得到这个局部核的 \(L^2\to L^6\) 估计，两个端点常数可分别取 \(\tfrac32(4\pi/3)^{1/3}\) 与 \(2\pi\)。这是本命题给出的一种充分指数范围。若删除核的空间截断，输入 \(h=1\) 已使远场积分发散，所用的 \(L^\infty\) 端点便不成立；全空间分数积分需要另行估计远场。

\subsection{极大算子应用}
\label{b2:subsec:170304}

\Needspace{6\baselineskip}
\begin{corollary}\label{b2:thm:17-maximal-strong}
设 \(d\) 为正整数，\(1<p\leq\infty\)。在 \(\mathbb R^d\) 上取 Lebesgue 测度，\(M\)、\(M_u\) 分别为中心球和包含观察点的球上绝对平均的上确界；\(M_{\mathcal D}\) 为标准半开二进网格
\(\mathcal D=\{\,2^{-j}(k+[0,1)^d):j\in\mathbb Z,\ k\in\mathbb Z^d\,\}\)
上包含观察点的立方体绝对平均的上确界。则三个算子都把 \(L^p(\mathbb R^d;\mathbb C)\) 映入自身，并满足
\[
 \|Mf\|_p+\|M_uf\|_p+\|M_{\mathcal D}f\|_p
 \leq C_{d,p}\|f\|_p.
\]
各算子在 \(L^\infty\) 上的范数均不超过一。
\end{corollary}
\begin{proof}
\cref{b2:thm:17-maximal-weak}及 \(M_u\leq2^dM\) 给出前两个算子的弱 \((1,1)\) 界，\cref{b2:thm:17-dyadic-level}给出二进版本。每个绝对平均都不超过输入的本性上确界，故三个算子均为 \(L^\infty\) 压缩。在 \(L^1+L^\infty\) 上，可积部分的极大函数几乎处处有限，有界部分的极大函数处处有界；次可加性因而使三个算子都具有所需定义域。应用\cref{b2:thm:17-marcinkiewicz}的 \(p_0=1,p_1=\infty\) 情形，即得有限 \(p\) 的结论。
\end{proof}

对中心极大算子，还可保留一个以后用于端点问题的精细分布界。设 \(d\geq1\)、\(1\leq p<\infty\)、\(f\in L^p(\mathbb R^d)\)、\(t>0\)。分解
\(f=f\cdot\mathbf1_{\{|f|>t/2\}}+f\cdot\mathbf1_{\{|f|\leq t/2\}}\)，
大值部分可积，小值部分的极大平均不超过 \(t/2\)。对大值部分用弱型估计，得到
\begin{equation}\label{b2:eq:17-maximal-tail-refined}
 |\{Mf>t\}|\leq\frac{2\cdot5^d}{t}
                   \int_{|f|>t/2}|f|\,\mathrm dx.
\end{equation}
当 \(p>1\) 时，保留截断再作高度积分，便有
\[
 \|Mf\|_p^p\leq \frac{5^dp\,2^p}{p-1}\|f\|_p^p.
\]
分母 \(p-1\) 来自 \(\int_0^{2|f(x)|}t^{p-2}\,\mathrm dt\)。

\ProseParagraphBreak

这一常数在 \(p\downarrow1\) 时的增长，可以用具体输入检验。在实直线上取 \(f=\mathbf1_{[0,1]}\)。当 \(x>2\) 时，半径为 \(x\) 的中心区间给出 \(Mf(x)\geq1/(2x)\)，所以
\[
 \|Mf\|_p^p\geq2^{-p}\int_2^\infty x^{-p}\,\mathrm dx
 =\frac{2^{1-2p}}{p-1},
 \quad \|f\|_p=1.
\]
因此算子范数至少为一个固定常数乘以 \((p-1)^{-1/p}\)。在 \(1<p\leq2\) 上，这与 \((p-1)^{-1}\) 可比：若 \(\delta=p-1\in(0,1]\)，则
\[
 \frac{\delta^{-1/p}}{\delta^{-1}}
 =\delta^{\delta/(1+\delta)}\in[e^{-1/e},1].
\]
结合上界，实直线上的中心极大算子范数在 \(p\downarrow1\) 时确有 \(1/(p-1)\) 的量级。精确最佳常数还依赖所选的极大算子，不由这两个界确定。

\ProseParagraphBreak

极大函数还能同时控制一整族正则化算子，进而说明它们在 Lebesgue 点恢复原函数。下面的核可取复值，不要求非负。

\Needspace{6\baselineskip}
\begin{proposition}\label{b2:prop:17-approximate-identity}
设 \(d\) 为正整数，\(\psi\in L^1(\mathbb R^d;\mathbb C)\)、\(\int_{\mathbb R^d}\psi=1\)。假设存在 \(\varepsilon>0\)、\(0<C<\infty\)，使
\[
 |\psi(x)|\leq C(1+|x|)^{-d-\varepsilon}
 \quad\text{几乎处处}.
\]
对 \(t>0\)，置 \(\psi_t(x)=t^{-d}\cdot\psi(x/t)\)。若 \(1\leq p<\infty\)、\(f\in L^p(\mathbb R^d;\mathbb C)\)，则 \(\psi_t*f\to f\) 于 \(L^p\)。对每个满足 \(Mf(x)<\infty\) 的点，所有 \(t>0\) 的卷积积分均绝对收敛，且
\[
 \sup_{t>0}|\psi_t*f(x)|\leq C'\cdot Mf(x),
\]
其中 \(M\) 为中心 Hardy--Littlewood 极大算子，\(C'\) 只依赖 \(d,\varepsilon,C\)。对任一有限值可测代表 \(f\)，在它的每个 Lebesgue 点 \(x\) 都有 \(\psi_t*f(x)\to f(x)\) 当 \(t\downarrow0\)。
\end{proposition}
\begin{proof}
先说明共同的极大函数控制。把积分区域分为 \(|y|<t\) 及环带 \(2^{j-1}t\leq|y|<2^jt\)、\(j\geq1\)。内球上 \(|\psi_t|\leq Ct^{-d}\)，相应积分不超过 \(C\omega_d\cdot Mf(x)\)。第 \(j\) 个环带上，核的上界为
\(C2^{d+\varepsilon}t^{-d}2^{-j(d+\varepsilon)}\)，而 \(|f(x-y)|\) 的积分不超过
\(\omega_d(2^jt)^d\cdot Mf(x)\)。所以
\[
 \int_{\mathbb R^d}|\psi_t(y)|\cdot|f(x-y)|\,\mathrm dy
 \leq C_{d,\varepsilon}C
       \left(1+\sum_{j\geq1}2^{-j\varepsilon}\right)\cdot Mf(x).
\]
右端与 \(t\) 无关，且几何级数收敛。这给出所述上确界估计及全部尺度的绝对收敛；由弱型或强型极大估计，满足 \(Mf(x)<\infty\) 的点组成满测集。

范数收敛可以回用\cref{b2:thm:03-approximate-identity}：缩放保持 \(\int\psi_t=1\)、\(\|\psi_t\|_1=\|\psi\|_1\)，且对每个 \(a>0\)，
\(\int_{|y|\geq a}|\psi_t(y)|\,\mathrm dy
=\int_{|z|\geq a/t}|\psi(z)|\,\mathrm dz\to0\)。
其证明在本处也可直接写成
\[
 \|\psi_t*f-f\|_p
 \leq\int_{\mathbb R^d}|\psi(z)|\cdot
             \|f(\,\cdot-tz)-f\|_p\,\mathrm dz.
\]
有限指数的平移连续性使每个固定 \(z\) 的范数差趋零，而共同可积控制为
\(2\|f\|_p\cdot|\psi(z)|\)；支配收敛给出结论。

固定 \(f\) 的一个 Lebesgue 点 \(x\)，令
\[
 A_x(r)=\frac1{|B(x,r)|}\int_{B(x,r)}|f(y)-f(x)|\,\mathrm dy.
\]
有 \(A_x(r)\to0\) 当 \(r\downarrow0\)。此外 \(S_x=\sup_{r>0}A_x(r)<\infty\)：先选 \(r_0>0\) 控制 \(r<r_0\) 的平均；对 \(r\geq r_0\)，Hölder 给出
\[
 A_x(r)\leq (\omega_dr^d)^{-1/p}\|f\|_p+|f(x)|
 \leq(\omega_dr_0^d)^{-1/p}\|f\|_p+|f(x)|.
\]
这也说明 \(Mf(x)\leq S_x+|f(x)|<\infty\)，从而该点的卷积积分已有定义。由 \(\int\psi_t=1\)，将前面的环带估计用于 \(|f(x-y)-f(x)|\)，得到
\[
 |\psi_t*f(x)-f(x)|
 \leq C_{d,\varepsilon}C
       \left(A_x(t)+\sum_{j\geq1}2^{-j\varepsilon}A_x(2^jt)\right).
\]
给定误差要求，先利用 \(A_x\leq S_x\) 与几何级数选出一个固定的尾部起点，使尾和对所有 \(t\) 都小；剩下只有有限个半径，每个都随 \(t\downarrow0\) 趋零。因此右端趋零。
\end{proof}

这里恢复的是所选代表在 Lebesgue 点的取值。\cref{b2:thm:17-differentiation}保证这些点几乎处处存在；一个 \(L^p\) 等价类本身并不指定任意一点的函数值。

% !TeX root = ../../Book2.tex
\section{Riesz–Thorin 插值}
\label{b2:sec:1704}

上一节把输入按大小分开，再对输出的高度积分。若算子本身是线性的，而且两个端点已经给出强范数估计，还可以让输入随一个复参数变化：在条带两侧分别适合两个端点，在条带内部指定的位置回到原输入。输出再与一族测试函数配对，就把算子估计变成一个复函数的边界估计。这样得到的常数是两个端点常数的几何平均。

\ProseParagraphBreak

这条路线需要保留复线性结构。极大算子的绝对值和上确界一般不随复参数全纯，不能直接放入本节的证明；Fourier 变换、固定核卷积等线性算子则可以。先证明实际需要的三线估计，再从竖边的范数要求推导测试函数的选择。

\ProseParagraphBreak

这里所用的最大模原理只涉及有界矩形 \(Q\)：若 \(H\) 在矩形内部全纯并连续到闭矩形，则内部的 \(|H|\) 不超过边界上的最大值。所需的基本复分析事实是全纯函数在内部光滑，且实部、虚部调和；写 \(H=u+\mathrm iv\)，Cauchy--Riemann 方程 \(u_x=v_y\)、\(u_y=-v_x\) 给出 \(\Delta u=v_{yx}-v_{xy}=0\)，虚部同理。因此对每个单位复数 \(c\)，可把本卷\cref{b2:cor:13-boundary-maximum}用于实值调和函数 \(\operatorname{Re}(cH)\)。给定内点 \(z\)，选择 \(c\) 使 \(cH(z)=|H(z)|\)，便得到
\[
 |H(z)|=\operatorname{Re}(cH(z))
 \leq\max_{\partial Q}\operatorname{Re}(cH)
 \leq\max_{\partial Q}|H|.
\]
若 \(H(z)=0\)，最后的模界直接成立。这里全纯指在复开域中处处存在复导数；光滑性与 Cauchy--Riemann 方程是本处使用的局部复分析基础。相同的调和最大值论证也适用于有界扇形等有界区域，只要函数在内部全纯并连续到整个闭域。闭矩形的紧性保证边界最大值存在；下面仍须先控制每个有限矩形的横边，才能令矩形高度趋于无穷，不能直接把这条有界域结论套到无限条带。

\subsection{三线引理与解析族}
\label{b2:subsec:170401}

下面采用\term[hadamard-sanxian-yinli]{Hadamard 三线引理}[Hadamard three-lines lemma]的下述版本。

\Needspace{6\baselineskip}
\begin{lemma}[Hadamard 三线估计]\label{b2:thm:17-three-lines}
设 \(M_0,M_1\geq0\)。复值函数 \(F\) 在开条带 \(0<\operatorname{Re}z<1\) 全纯，在闭条带 \(0\leq\operatorname{Re}z\leq1\) 连续且有界，并对每个 \(t\in\mathbb R\) 满足 \( |F(\mathrm it)|\leq M_0\)、\( |F(1+\mathrm it)|\leq M_1\)。则对每个 \(0<\theta<1\)，有
\[
 |F(\theta)|\leq M_0^{1-\theta}M_1^\theta.
\]
\end{lemma}
\begin{proof}
先设 \(M_0,M_1>0\)，令 \(G(z)=F(z)M_0^{z-1}M_1^{-z}\)，幂以正实数的实对数定义。若 \(z=s+\mathrm it\)，则归一化因子的模为 \(M_0^{s-1}M_1^{-s}\)，与 \(t\) 无关；所以 \(G\) 在闭条带有界，两条竖边的模不超过一。

竖边界尚未控制有限矩形的横边。为此固定 \(\varepsilon>0\)，乘上 \(e^{\varepsilon(z^2-1)}\)。它在 \(s=0,1\) 时的模不超过一，在 \(t=\pm R\)、\(0\leq s\leq1\) 时的模不超过 \(e^{-\varepsilon R^2}\)。若 \(|G|\leq B\)，只须选 \(R\) 足够大使 \(B e^{-\varepsilon R^2}\leq1\)，四条边就都受一控制。矩形最大模原理给出
\[
 |G(\theta)|e^{\varepsilon(\theta^2-1)}\leq1,
 \quad
 |F(\theta)|\leq M_0^{1-\theta}M_1^\theta
 e^{\varepsilon(1-\theta^2)}.
\]
令 \(\varepsilon\downarrow0\) 得到结论。若某个 \(M_i=0\)，先以正数 \(M_i+\delta\) 作为两侧较大的界，再令 \(\delta\downarrow0\)。这同时处理了一侧或两侧恒为零的情形。
\end{proof}

三线估计要求的对象是一个在整个条带中有界的标量全纯函数，不只是两条竖边上的范数公式。下面取有限个简单函数的组合，既能逐项核对全纯性，也能直接证明内部有界。

\ProseParagraphBreak

设 \(U\subset\mathbb C\) 开，\(T_z\) 是定义在同一复测试向量空间上、\(z\in U\) 的复线性算子族，并固定可配对的测试输出空间。一个\term[jiexi-zu]{解析族}[analytic family]是复参数下的算子族，其对测试输入、测试输出的矩阵元为全纯函数。使用三线估计时必须检查边界范数和条带内增长；“公式含有 \(z\)”不是解析族的充分条件。下面先用固定算子和变化的测试函数，避免讨论一般算子的定义域变化。

\subsection{Riesz–Thorin 定理}
\label{b2:subsec:170402}

两个端点必须作用于同一个输入，并在交空间上给出同一输出。例如，一个端点从积分定义算子，另一个从完备化定义算子，就要先核对这两种定义在交空间上一致。此时可令 \(T(f_0+f_1)=T_0f_0+T_1f_1\)：交空间上的一致性保证结果与分解无关。下面直接把这项相容性写成和空间上的一个复线性算子。

\ProseParagraphBreak

本节把\cref{b2:subsec:04-complex-interpolation}中出现的 Riesz--Thorin 方法推广到一般线性算子，允许输入、输出的两个端点指数分别变化。约定 \(1/\infty=0\)，\(r'\) 表示满足 \(1/r+1/r'=1\) 的 Hölder 共轭指数。

\Needspace{6\baselineskip}
\begin{theorem}[Riesz--Thorin]\label{b2:thm:17-riesz-thorin}
设 \((X,\mathcal A,\mu)\)、\((Y,\mathcal B,\nu)\) 为 \(\sigma\)-有限测度空间，所有 \(L^r\) 空间均取复值函数的几乎处处等价类。设 \(1\leq p_i,q_i\leq\infty\)、\(A_i\geq0\)、\(i=0,1\)，且 \(T\) 是从 \(L^{p_0}(X,\mu)+L^{p_1}(X,\mu)\) 到 \(Y\) 上几乎处处有限的可测函数等价类空间的复线性算子，满足
\[
 \|Tf\|_{L^{q_i}(Y,\nu)}\leq A_i\|f\|_{L^{p_i}(X,\mu)}
 \quad\bigl(f\in L^{p_i}(X,\mu),\ i=0,1\bigr).
\]
给定 \(0<\theta<1\)，以 \(1/\infty=0\) 定义 \(p,q\)：
\[
 \dfrac1p=\dfrac{1-\theta}{p_0}+\dfrac\theta{p_1},\quad
 \dfrac1q=\dfrac{1-\theta}{q_0}+\dfrac\theta{q_1},
\]
则每个 \(f\in L^p(X,\mu)\) 都属于上述和空间，且
\[
 \|Tf\|_{L^q(Y,\nu)}
 \leq A_0^{1-\theta}A_1^\theta\|f\|_{L^p(X,\mu)}.
\]
\end{theorem}
\begin{proof}
先证明有限测度支集上的简单输入的估计。零输入无需处理；将输入 \(f\) 与有限测度支集上的简单测试函数 \(g\) 归一化为
\(\|f\|_p=\|g\|_{q'}=1\)。先设 \(p,q'<\infty\)。希望构造 \(f_z\)，使它在 \(z=\theta\) 等于 \(f\)，在 \(\operatorname{Re}z=i\)、\(i=0,1\) 时的 \(L^{p_i}\) 范数等于一。保留 \(f\) 的复相位，试取
\[
 f_z=\frac f{|f|}\cdot|f|^{\alpha(z)}.
\]
若 \(p_i<\infty\)，竖边的归一化要求
\(p_i\operatorname{Re}\alpha(i+\mathrm it)=p\)，这样积分恰好仍是 \(\int|f|^p=1\)。若 \(p_i=\infty\)，则要求实部为零，使支集上模恒为一。这两种要求统一为 \(\alpha(i)=p/p_i\)。取连接这两个值的仿射函数，就得到
\[
 \alpha(z)=p\left(\frac{1-z}{p_0}+\frac z{p_1}\right),
 \quad \alpha(\theta)=1.
\]
输出使用 Hölder 配对，所以应对 \(g\) 插值的是共轭指数。由于
\[
 \frac1{q'}=\frac{1-\theta}{q_0'}+\frac\theta{q_1'},
\]
相同的选择给出
\[
 \beta(z)=q'\left(\frac{1-z}{q_0'}+\frac z{q_1'}\right),
 \quad
 g_z=\frac g{|g|}\cdot|g|^{\beta(z)},
 \quad g_\theta=g.
\]
所有复幂都以非零函数值的模的实对数定义；原函数为零的位置统一赋零。即使某侧实部为零，也不把支集外的零改成一。在竖边 \(i=0,1\) 上，以上计算给出
\[
 \|f_{i+\mathrm it}\|_{p_i}
 =\|g_{i+\mathrm it}\|_{q_i'}=1.
\]
这里含无穷端点的等式，是相应函数在其非零支集上模恒为一，而非使用任何无穷次幂。

还要补足中间指数本身为无穷的构造。若 \(p=\infty\)，倒数插值式迫使 \(p_0=p_1=\infty\)。此时对任意归一化的 \(f\in L^\infty(X)\) 直接取 \(f_z=f\)，不要求它是简单函数或具有有限测度支集。若 \(q'=\infty\)，则 \(q_0'=q_1'=\infty\)，即两个输出端点均为 \(L^1\)；这时取 \(g_z=g\)，其中 \(g\) 仍为归一化的有限测度支集简单函数。其余一侧仍按上面的有限指数公式构造。这样所有情形都满足竖边范数不超过一，并在 \(\theta\) 处恢复原函数。

定义双线性配对
\[
 F(z)=\int_Y (Tf_z)(y)\cdot g_z(y)\,\mathrm d\nu(y).
\]
这里不对 \(g_z\) 取复共轭，否则参数一般会变成反全纯。对偶测试所需的共轭相位可以预先放进 \(g\) 中。若 \(f=\sum_j a_j\cdot\mathbf1_{E_j}\)、\(g=\sum_k b_k\cdot\mathbf1_{H_k}\)，则 \(F\) 是有限个整函数系数的乘积与固定数
\(\int_{H_k}T\mathbf1_{E_j}\,\mathrm d\nu\) 的线性组合。\(E_j,H_k\) 测度有限，端点估计及 Hölder 保证这些积分绝对收敛。若 \(p=\infty\)，只把其中的 \(T\mathbf1_{E_j}\) 换为固定的 \(Tf\)，同样由一个输出端点保证局部可积。因此 \(F\) 全纯并连续到闭条带，无需未经控制地交换求导和积分。

在闭条带上，复幂系数的模只依赖 \(\operatorname{Re}z\in[0,1]\)，因而各自有界；有限和使 \(F\) 在整个闭条带有界。两条竖边则满足
\[
 |F(i+\mathrm it)|
 \leq\|Tf_{i+\mathrm it}\|_{q_i}
 \cdot\|g_{i+\mathrm it}\|_{q_i'}
 \leq A_i.
\]
由\cref{b2:thm:17-three-lines}，在 \(z=\theta\) 得到
\[
 \left|\int_Y (Tf)\cdot g\,\mathrm d\nu\right|
 \leq A_0^{1-\theta}A_1^\theta.
\]
恢复两个归一化因子后，右侧再乘 \(\|f\|_p\cdot\|g\|_{q'}\)。

现在对所有这些 \(g\) 取上确界。对于在每个有限测度集上可积的可测函数 \(h\)，所用的范数刻画是
\[
 \|h\|_q
 =\sup_{\substack{g\text{ 为有限测度支集简单函数}\\\|g\|_{q'}\leq1}}
 \left|\int_Y h\cdot g\,\mathrm d\nu\right|.
\]
这个等式允许两边取无穷。对 \(1<q<\infty\)，先在有限测度集及 \(|h|\) 有界的部分上，取共轭相位 \(\overline h/|h|\) 乘 \(|h|^{q-1}\)，再归一化、简单逼近并穷竭；此过程也检验了原先未知的 \(L^q\) 可积性。对 \(q=1\)，测试函数改为该共轭相位在有限测度集上的截断；对 \(q=\infty\)，在 \(|h|\) 接近本性上确界的有限正测度子集上，使用该共轭相位除以集合测度的测试函数。所有相位在 \(h=0\) 处赋零。\(\sigma\)-有限性保证最后这样的子集能够选出，也允许用有限测度集穷竭 \(Y\)。当前的 \(Tf\) 在两个输出端点中，Hölder 保证它在每个有限测度集上可积，故上述测试适用。由此得到简单输入的估计；\(p=\infty\) 时，前面已经对任意有界输入证明了结论，无需再用稠密性。

最后设 \(p<\infty\)。\(\sigma\)-有限性允许先截取有限测度支集，再截断幅度并作简单逼近，得到 \(f_n\rightarrow f\) 于 \(L^p(X)\)。已证估计使 \(Tf_n\) 在 \(L^q(Y)\) 中 Cauchy，即使 \(q=\infty\) 也可用该空间的完备性，记极限为 \(v\)。这定义了唯一的有界线性延拓，但仍须核对 \(v\) 等于题设和空间中原有的 \(Tf\)。

为此将两个源端点按大小记为 \(r_0\leq p\leq r_1\)。对非零 \(h\in L^p\)，令 \(a=\|h\|_p\)，以 \(|h|>a\) 与 \(|h|\leq a\) 分解为 \(h_0+h_1\)。幂次比较给出
\[
 \|h_0\|_{r_0}\leq\|h\|_p,
 \quad \|h_1\|_{r_1}\leq\|h\|_p;
\]
若 \(r_1=\infty\)，第二项直接由 \(|h_1|\leq a\) 得到。于是 \(L^p\) 连续嵌入原来的端点和空间。对 \(h=f_n-f\) 使用这个分解，端点估计使 \(T h_0,T h_1\) 在各自的输出端点范数中趋零。Chebyshev 不等式说明它们在每个有限测度目标集上依测度趋零；无穷输出端点则直接给一致的本性上界趋零。因此 \(Tf_n\rightarrow Tf\) 在这些集合上依测度成立。另一方面 \(Tf_n\rightarrow v\) 的 \(L^q\) 收敛也蕴含同样的依测度收敛。极限唯一性和 \(Y\) 的可数穷竭给出 \(v=Tf\) 几乎处处，从而原算子本身满足所述估计。
\end{proof}

\begin{table}[htbp]
\centering
\caption[解析测试族的指数选择]{解析测试族的指数选择。有限中间指数时，复幂的实部使竖边范数归一，插值点的幂则为一；约定 \(1/\infty=0\)。中间指数 \(p\) 或 \(q'\) 本身为无穷时，相应的两个测试指数 \(p_i\) 或 \(q_i'\) 也均为无穷，应固定原函数，不能计算含 \(\infty\) 的复幂。}
\label{b2:tab:17-analytic-test-exponents}
\begin{tabularx}{\linewidth}{@{}p{9.4em}XXX@{}}
\toprule
测试族 & 左竖边的实部 & 右竖边的实部 & 插值点的取值\\
\midrule
\(p<\infty\) 的 \(f_z\) & \(\operatorname{Re}\alpha=p/p_0\) & \(\operatorname{Re}\alpha=p/p_1\) & \(\alpha(\theta)=1\)\\
\(q'<\infty\) 的 \(g_z\) & \(\operatorname{Re}\beta=q'/q_0'\) & \(\operatorname{Re}\beta=q'/q_1'\) & \(\beta(\theta)=1\)\\
\(p=\infty\) 的 \(f_z\) & 固定 \(f\) & 固定 \(f\) & \(f_\theta=f\)\\
\(q'=\infty\) 的 \(g_z\) & 固定 \(g\) & 固定 \(g\) & \(g_\theta=g\)\\
\bottomrule
\end{tabularx}
\end{table}
\cref{b2:tab:17-analytic-test-exponents}把两个作用分开：边界实部决定可用的端点范数，虚部只改变相位；插值点的幂等于一，才使最终估计针对原来的输入与测试函数。

\ProseParagraphBreak

解析算子族 \(T_z\) 的使用也应归结为同样的标量检验：对所选的变化测试函数，\(F(z)=\int(T_zf_z)\cdot g_z\) 必须在条带内部全纯、连续到边界。单独知道固定测试函数的矩阵元全纯时，还须说明变化测试族是有限解析组合，或给出相应的积分控制。若 \(F\) 在整个条带至多按 \(C_{f,g}e^{c|\operatorname{Im}z|}\) 增长，而两条竖边仍有常数界 \(A_i\)，对任意 \(\delta>0\) 乘上 \(e^{\delta(z-\theta)^2}\) 后便可用有界版本；其竖边界增加的因子至多是 \(e^{\delta(i-\theta)^2}\)。令 \(\delta\downarrow0\)，仍得精确的几何平均常数。

\ProseParagraphBreak

后文的某些族在竖边上也会增长。若对所有归一化测试函数有共同的 \(c\geq0\)，并且
\[
 |F(s+\mathrm it)|\leq C_{f,g}e^{ct^2}
 \quad(0\leq s\leq1),
 \quad |F(i+\mathrm it)|\leq A_i e^{ct^2}\quad(i=0,1),
\]
则取 \(\delta>c\) 后，同一乘子抑制了内部增长。应用三线估计得到
\[
 |F(\theta)|
 \leq e^{\delta\theta(1-\theta)}A_0^{1-\theta}A_1^\theta.
\]
因为 \(c\) 对测试函数一致，附加常数也一致；这里不能再令 \(\delta\) 趋零。曲面限制中的解析族允许这种 Gaussian 型增长。增长条件及其常数依赖必须在具体应用中验证，不能只引用有界版本后略去它们。

\ProseParagraphBreak

先看一个能直接算出尺度的应用。设 \(d\) 为正整数、\(\kappa,t>0\)，取热核
\[
 H_{\kappa,t}(x)=(4\pi\kappa t)^{-d/2}
 e^{-|x|^2/(4\kappa t)},
 \quad T_t f=H_{\kappa,t}*f.
\]
固定 \(1\leq p<\infty\)。核的积分为一，所以 Young 不等式给出 \(L^p\rightarrow L^p\) 的常数一；Hölder 给出 \(L^p\rightarrow L^\infty\) 的常数 \(\|H_{\kappa,t}\|_{p'}\)。这两个端点使用同一个积分卷积，故相容。对 \(p<q<\infty\)，输入指数保持为 \(p\)，目标指数满足 \(1/q=(1-\theta)/p\)，即 \(\theta=1-p/q\)。Gaussian 积分算出
\[
 \|H_{\kappa,t}\|_r
 =r^{-d/(2r)}(4\pi\kappa t)^{-\frac d2(1-1/r)}
 \quad(1\leq r<\infty),
 \quad \|H_{\kappa,t}\|_\infty=(4\pi\kappa t)^{-d/2}.
\]
因此
\[
 \|T_t f\|_q
 \leq C_{d,p,q}(4\pi\kappa t)^{-\frac d2(1/p-1/q)}\|f\|_p,
 \quad
 C_{d,p,q}=(p')^{-\frac{d(1-p/q)}{2p'}}.
\]
当 \(p=1\) 时该常数按 \(p'=\infty\) 的极限取一。\(q=p\) 与 \(q=\infty\) 分别由刚才两个端点直接得到；\(p=q=\infty\) 也由核的积分为一得到压缩。譬如 \(p=2,q=4\) 时，
\[
 \|T_t f\|_4\leq 2^{-d/8}(4\pi\kappa t)^{-d/8}\|f\|_2.
\]
正时间带来更高的可积性，而常数在 \(t\downarrow0\) 时按明确的幂发散。插值没有丢失尺度：两个端点中 \(\kappa t\) 的幂，按同一个 \(\theta\) 组合。这也解释了为什么应先记录尺度依赖，再对许多尺度求和。

\subsection{Hausdorff–Young 再论}
\label{b2:subsec:170403}

\Needspace{6\baselineskip}
\begin{corollary}[Hausdorff--Young]\label{b2:thm:17-hausdorff-young}
设 \(d\) 为正整数，\(1\leq p\leq2\)，\(p'\) 为 Hölder 共轭指数（\(p=1\) 时 \(p'=\infty\)）。对 \(f\in L^p(\mathbb R^d;\mathbb C)\)，令 \(\widehat f\) 为使用核 \(e^{-2\pi\mathrm ix\cdot\xi}\) 的可积 Fourier 变换或其从有限测度支集简单函数延拓所得的变换。则
\[
 \|\widehat f\|_{p'}\leq\|f\|_p.
\]
\end{corollary}
\begin{proof}
可积 Fourier 变换满足 \(L^1\rightarrow L^\infty\) 的范数界一；\cref{b2:thm:04-plancherel-extension}给出 \(L^2\rightarrow L^2\) 的等距延拓，并明确说明它在 \(L^1\cap L^2\) 上与可积变换一致。因此二者给出同一个和空间算子。对 \(1<p<2\)，取
\[
 \theta=2-\frac2p,
 \quad \frac1p=1-\frac\theta2,
 \quad \frac1{p'}=\frac\theta2.
\]
将端点 \((p_0,q_0)=(1,\infty)\)、\((p_1,q_1)=(2,2)\) 及 \(A_0=A_1=1\) 代入\cref{b2:thm:17-riesz-thorin}，即得所需估计。两个端点由原有结论直接处理。对有限测度支集简单函数的 \(L^p\) 逼近产生唯一连续延拓；对一般 \(L^p\) 输入，这个结论不要求普通 Fourier 积分逐点绝对收敛。
\end{proof}
常数一来自两个端点的相同上界。Gaussian 的可计算范数已经提示最优常数问题比有界性更细。逆变换具有相同的两个端点，因而同样得到 \(L^p\) 到 \(L^{p'}\) 的估计。

\ProseParagraphBreak

指数上限也能用本章可完成的有限计算检验。设 \(d\) 为正整数、\(2<p\leq\infty\)，置 \(r=p'\in[1,2)\)。选非负、非零的 \(\psi\in C_c^\infty((0,1/4)^d)\)，并对符号 \(\varepsilon_k\in\{-1,1\}\) 令
\[
 f_N(x)=\sum_{k=1}^N\varepsilon_k\cdot\psi(x-ke_1),
 \quad
 \widehat f_N(\xi)=\widehat\psi(\xi)\cdot
 \sum_{k=1}^N\varepsilon_k e^{-2\pi\mathrm ik\xi_1}.
\]
平移后的支集不交，所以
\(\|f_N\|_p=N^{1/p}\|\psi\|_p\)，包括 \(p=\infty\) 时的约定 \(N^{1/p}=1\)。若符号全部为正，只在 \(|\xi_1|\lesssim1/N\) 的窄峰上计算，所得 \(L^r\) 下界是 \(N^{1-1/r}=N^{1/p}\)，与输入同阶，不能据此排除有界性。需要让变换在固定大小的区域上保留约 \(\sqrt N\) 的平均量级。

\ProseParagraphBreak

这里的符号平均只是在全部 \(2^N\) 种选择上取算术平均，记作 \(\mathbb E\)；事件的概率 \(\mathbb P\) 就是相应选择所占的比例。对固定 \(\xi\)，令 \(S_N=\sum_k\varepsilon_k e^{-2\pi\mathrm ik\xi_1}\)。展开乘积时，含有某个符号的奇数次幂的项平均为零；在四次乘积中，留下的下标至多有三种两两配对方式。因此
\[
 \mathbb E|S_N|^2=N,
 \quad \mathbb E|S_N|^4\leq3N^2.
\]
设 \(B=\{\,|S_N|^2\geq N/2\,\}\)。其补集上二次量小于 \(N/2\)，故 Cauchy--Schwarz 给出
\[
 \frac N2\leq\mathbb E\bigl(|S_N|^2\cdot\mathbf1_B\bigr)
 \leq\bigl(\mathbb E|S_N|^4\bigr)^{1/2}\cdot\mathbb P(B)^{1/2}.
\]
于是 \(\mathbb P(B)\geq1/12\)，并有
\[
 \mathbb E|S_N|^r\geq\frac1{12}(N/2)^{r/2}.
\]
由于 \(\widehat\psi(0)=\int\psi>0\) 且变换连续，存在固定的有限正测度频率盒 \(J\) 与 \(a>0\)，使 \(|\widehat\psi|\geq a\) 于 \(J\)。对 \(\xi\in J\) 积分并与有限符号平均交换，得到
\[
 \mathbb E\int_J|\widehat f_N(\xi)|^r\,\mathrm d\xi
 \geq\frac{a^r\cdot|J|}{12}(N/2)^{r/2}.
\]
至少有一组符号使 \(\|\widehat f_N\|_r\geq c_{\psi,r}\sqrt N\)。因此
\[
 \frac{\|\widehat f_N\|_{p'}}{\|f_N\|_p}
 \geq c_{\psi,p}N^{1/2-1/p}\longrightarrow\infty.
\]
这些输入本身光滑且紧支撑，已经排除了 \(p>2\) 时任何与输入无关的 \(L^p\rightarrow L^{p'}\) 常数。符号的作用是让固定频率区域上的平均大小保持为平方根量级，而不是只制造越来越窄的同号峰。

\subsection{插值方法与 \texorpdfstring{\(K\)}{K} 泛函}
\label{b2:subsec:170404}

Marcinkiewicz 方法使用大小分解和分布函数，能处理极大算子这样的次线性对象；Riesz--Thorin 方法保留复线性结构，从两个强型端点得到几何平均常数。对奇异积分，常见路线是先用弱 \((1,1)\) 与强 \((2,2)\) 得到 \(1<p<2\)，再由伴随算子的估计处理另一半指数。选用哪种方法，取决于已知端点的实际强弱，不能把弱型范数直接当作解析插值中的强型范数。

\ProseParagraphBreak

大小分解也提示另一种问题：不预先规定截断高度，而把一个输入分配给两个空间，能达到的最小总代价是什么？设 \(X_0,X_1\) 是同一实或复标量域上的 Banach 空间，连续嵌入同一个 Hausdorff 线性拓扑空间，使 \(f_0+f_1\) 有共同的意义。对 \(f\in X_0+X_1\) 及 \(t>0\)，定义其 \term[k-fanhan]{\(K\) 泛函}[K-functional]
\[
 K(t,f;X_0,X_1)\coloneqq
 \inf_{\substack{f=f_0+f_1\\f_i\in X_i}}
 \bigl(\|f_0\|_{X_0}+t\|f_1\|_{X_1}\bigr).
\]
参数 \(t\) 调整第二个空间的代价。\(t\) 小时可以把更多部分交给 \(X_1\)，\(t\) 大时则倾向于减少 \(X_1\) 的范数。因此下确界中的分解必须允许随 \(t\) 改变。对 \(L^1,L^\infty\)，有限的 \(L^\infty\) 代价对应削平函数的大值，剩余高度由 \(L^1\) 支付；这个分配可以逐点调整幅度，不要求把某个水平集切成任意指定测度的子集。

\Needspace{6\baselineskip}
\begin{proposition}\label{b2:prop:17-k-functional}
设 \((X,\mathcal A,\mu)\) 为 \(\sigma\)-有限测度空间，\(f\in L^1(X,\mu)+L^\infty(X,\mu)\) 为实或复值函数。令 \(d_f(\lambda)=\mu\{|f|>\lambda\}\)、\(\lambda\geq0\)，并定义递减重排 \(f^*(s)=\inf\{\lambda\geq0:d_f(\lambda)\leq s\}\)、\(s>0\)，空集的下确界取无穷。对每个 \(t>0\)，令 \(K(t,f;L^1,L^\infty)=\inf_{f=u+v}(\|u\|_1+t\|v\|_\infty)\)，其中 \(u\in L^1(X,\mu)\)、\(v\in L^\infty(X,\mu)\)。则
\[
 K(t,f;L^1,L^\infty)=\int_0^t f^*(s)\,\mathrm ds.
\]
另设 \((X_0,X_1)\)、\((Y_0,Y_1)\) 为同一实或复标量域上的相容 Banach 空间对（每对连续嵌入同一个 Hausdorff 线性拓扑空间），\(T:X_0+X_1\rightarrow Y_0+Y_1\) 是线性算子。若 \(T:X_i\rightarrow Y_i\) 的范数分别至多 \(A_i>0\)、\(i=0,1\)，则对每个 \(f\in X_0+X_1\) 和 \(t>0\)，有
\[
 K(t,Tf;Y_0,Y_1)\leq A_0\cdot K(tA_1/A_0,f;X_0,X_1).
\]
\end{proposition}
\begin{proof}
先把重排积分改写为分布函数积分。对 \(s>0\)、\(\lambda\geq0\)，若 \(d_f(\lambda)\leq s\)，则 \(\lambda\) 已属于定义下确界的集合，故 \(f^*(s)\leq\lambda\)。若 \(d_f(\lambda)>s\)，右连续性保证某个 \(\delta>0\) 满足 \(d_f(\lambda+\delta)>s\)，于是 \(f^*(s)\geq\lambda+\delta>\lambda\)。因此 \(f^*(s)>\lambda\) 当且仅当 \(s<d_f(\lambda)\)。对 \(s\in(0,t)\) 使用层蛋糕公式及 Tonelli 定理，得到
\[
 J_t\coloneqq\int_0^t f^*(s)\,\mathrm ds
 =\int_0^\infty\min\{t,d_f(\lambda)\}\,\mathrm d\lambda.
\]
这个恒等式暂时允许两边为无穷；当某个 \(d_f(\lambda)=\infty\) 时，被积函数的相应值就是有限数 \(t\)。因此不必假定各个超水平集合都有有限测度。

取任意分解 \(f=u+v\)，其中 \(u\in L^1\)、\(v\in L^\infty\)，并令 \(b=\|v\|_\infty\)。在 \(\lambda>b\) 时，\(|f|\leq|u|+b\) 几乎处处蕴含
\(d_f(\lambda)\leq d_u(\lambda-b)\)。将高度积分在 \(b\) 处分开，得到
\[
 J_t\leq tb+\int_b^\infty d_u(\lambda-b)\,\mathrm d\lambda
 =t\|v\|_\infty+\|u\|_1.
\]
最后一步是 \(u\) 的层蛋糕公式。至少存在一个这样的分解，所以 \(J_t<\infty\)；再对所有分解取下确界，得 \(J_t\leq K(t,f;L^1,L^\infty)\)。整个比较通过高度分布完成，因此也适用于带原子的测度空间。

反向令 \(a=f^*(t)\)。重排单调不增，所以 \(ta\leq J_t<\infty\)，从而 \(a\) 是有限非负数。由广义逆定义及分布函数的右连续性，\(d_f(\lambda)>t\) 对 \(0\leq\lambda<a\) 成立，而 \(d_f(\lambda)\leq t\) 对 \(\lambda\geq a\) 成立。因此
\[
 J_t=ta+\int_a^\infty d_f(\lambda)\,\mathrm d\lambda.
\]
当 \(a=0\) 时，第一段高度区间为空，等式仍成立。现在以保持相位的截断定义
\[
 f_1(x)=
 \begin{cases}f(x),&|f(x)|\leq a,\\a f(x)/|f(x)|,&|f(x)|>a,\end{cases}
 \quad f_0=f-f_1.
\]
这里 \(f_1\) 满足 \(\|f_1\|_\infty\leq a\)，且 \(|f_0|=(|f|-a)_+\)。再次由层蛋糕公式，
\[
 \|f_0\|_1=\int_a^\infty d_f(\lambda)\,\mathrm d\lambda
 =J_t-ta<\infty.
\]
因此这是一个合法的 \(L^1+L^\infty\) 分解，其代价至多为 \(J_t-ta+ta=J_t\)。两边合起来即得所需等式。水平集即使是不可分割的原子，仍可对该原子上的函数值作幅度截断，因而不影响这个分解。

对于算子估计，每个分解产生 \(Tf=Tf_0+Tf_1\)，代价至多 \(A_0\|f_0\|_{X_0}+tA_1\|f_1\|_{X_1}\)。对全部分解取下确界即可。
\end{proof}

例如，取单点空间 \(X=\{x_0\}\)，令 \(\mu(\{x_0\})=1\)、\(f(x_0)=1\)、\(t=1/2\)。它的递减重排在 \(0<s<1\) 上恒为一，故 \(K(1/2,f;L^1,L^\infty)=1/2\)。直接取 \(f_1=f\)、\(f_0=0\) 就达到这个代价。另一方面，满足 \(\mu(E)\leq1/2\) 的集合只有空集，因此
\[
 \sup_{\mu(E)\leq1/2}\int_E|f|\,\mathrm d\mu=0.
\]
原子不能切成半个集合，但可以把原子上的函数值截断。这个例子说明，上述集合积分上确界不能在一般测度空间中代替 \(K\) 泛函，重排公式仍然成立则来自幅度截断。

\ProseParagraphBreak

对相容 Banach 空间对 \((X_0,X_1)\)、\(0<\theta<1\)、\(1\leq q\leq\infty\)，将所有尺度的分解代价合起来，可定义实插值量
\[
 \|f\|_{\theta,q;K}\coloneqq
 \left(\int_0^\infty
 \bigl[t^{-\theta}\cdot K(t,f;X_0,X_1)\bigr]^q
 \frac{\mathrm dt}{t}\right)^{1/q}
 \quad(f\in X_0+X_1),
\]
当 \(q=\infty\) 时取 \(\sup_{t>0}t^{-\theta}\cdot K(t,f;X_0,X_1)\)。乘法测度 \(\mathrm dt/t\) 使不同数量级的尺度具有同样权重；因子 \(t^{-\theta}\) 则决定在两个端点之间的位置。

\ProseParagraphBreak

对同一 \(\sigma\)-有限测度空间上的 \(L^1,L^\infty\)，令 \(1/p=1-\theta\)，故 \(1<p<\infty\)、\(\theta=1/p'\)。对 \(f\in L^1+L^\infty\)，上一命题给出的公式使实插值量与\secref{b2:sec:1702}中定义的 Lorentz 量满足
\[
 \|f\|_{L^{p,q}}\leq\|f\|_{\theta,q;K}
 \leq p'\|f\|_{L^{p,q}}.
\]
先看左侧。由于 \(f^*\) 单调不增，
\[
 K(t,f)=\int_0^t f^*(s)\,\mathrm ds
 \geq t\cdot f^*(t),
 \quad t^{-1/p'}\cdot K(t,f)\geq t^{1/p}f^*(t).
\]
对 \(t\) 取 \(L^q(\mathrm dt/t)\) 范数便得所需下界。反向令 \(h(s)=s^{1/p}f^*(s)\)，换元 \(s=tu\) 得到
\[
 t^{-1/p'}\cdot K(t,f)
 =\int_0^1u^{-1/p}\cdot h(tu)\,\mathrm du.
\]
在 \(\mathrm dt/t\) 中，固定 \(u>0\) 的缩放 \(t\mapsto tu\) 保持范数。因此有限 \(q\geq1\) 时，Minkowski 给出
\[
 \|f\|_{\theta,q;K}
 \leq\int_0^1u^{-1/p}\,\mathrm du\cdot
 \|h\|_{L^q(\mathrm dt/t)}
 =p'\|f\|_{L^{p,q}}.
\]
\(q=\infty\) 时把上确界移入积分得到同一界。这里 \(p>1\) 保证 \(\int_0^1u^{-1/p}\,\mathrm du<\infty\)，\(q\geq1\) 则保证所用 Minkowski 有效。不能只保留形式相同的重排表达而删除这两个条件。

\ProseParagraphBreak

上一命题的算子估计又给出相容 Banach 空间对之间的结论。若复线性或实线性算子 \(T:X_0+X_1\rightarrow Y_0+Y_1\) 的两个端点范数分别至多 \(A_0,A_1>0\)，则
\[
 K(t,Tf;Y_0,Y_1)
 \leq A_0\cdot K(tA_1/A_0,f;X_0,X_1).
\]
在实插值量中令 \(u=tA_1/A_0\)，有
\[
 t^{-\theta}A_0\cdot K(tA_1/A_0,f)
 =A_0^{1-\theta}A_1^\theta u^{-\theta}\cdot K(u,f).
\]
由于 \(\mathrm dt/t=\mathrm du/u\)，对有限 \(q\) 积分，或在 \(q=\infty\) 时取上确界，得到
\[
 \|Tf\|_{\theta,q;K(Y_0,Y_1)}
 \leq A_0^{1-\theta}A_1^\theta
 \|f\|_{\theta,q;K(X_0,X_1)}.
\]
这里的下标仅说明分别用哪一对空间计算 \(K\) 泛函。端点常数的总次数为一，与算子缩放 \(T\mapsto cT\) 时范数乘 \(|c|\) 相符。这一计算解释了几何平均常数为什么在两种插值方法中都会出现：复方法来自三线估计，实方法来自分解代价的尺度换元。

\ProseParagraphBreak

实方法还记录 Lorentz 的第二指数，但上述 Banach 空间对估计并不自动给出任意两个弱型端点之间的最佳 Lorentz 范围。弱端点首先提供的是分布函数控制，能否使用等价范数、输入如何分解以及目标是否完备，都须按具体指数检查。固定的可积核卷积可以使用强型端点和解析插值；临界弱核常更适合大小分解及集合积分；需要区分整个高度分布时，\(K\) 泛函又提供了相应的量。

\ProseParagraphBreak

这些方法也都不能仅靠开放区间的结论恢复端点。极大函数在每个 \(p>1\) 上有界，但单位球示性函数的像有不可积尾部；当估计常数在 \(p\downarrow1\) 时发散，便没有可用于取极限的统一控制。类似地，Fourier 的两个端点给出中间强型界，并不决定最优常数、核的抵消或主值存在性。使用插值之后，仍应按当前问题分别证明这些额外性质。


\begin{chaptersummary}
极大函数的弱型估计控制的是“大于某个高度的点有多少”。单独一个弱端点仍可能产生对数发散；Marcinkiewicz 方法让输入截断随输出高度改变，再交换积分次序，才得到严格位于两个端点之间的强型估计。二进分解还提供坏块的平均上下界，为下一章利用零平均抵消准备了空间分块。

\ProseParagraphBreak

Riesz--Thorin 方法则把端点指数编码在解析测试函数中，由三线引理得到范数常数的几何平均。使用时先核实两端是同一个兼容线性算子，并区分有限指数的稠密延拓与无穷端点的直接论证。Lorentz 空间进一步区分临界幂次下的对数可积性，\(K\) 泛函把分解代价写成重排积分。后面选择哪种方法，取决于手中掌握的是弱型还是强型、算子是否线性，以及是否需要保留尺度参数和第二指数。
\end{chaptersummary}

\BookExercises

% 自拟：围绕本章已经建立的工具，改变条件或组合方法。
\Needspace{6\baselineskip}
\begin{exercise}\label{b2:ex:17-new-1}
在带 Lebesgue 测度的 \(\mathbb R\) 上令 \(f=\mathbf1_{(0,1)}\)。求 \(x>1\) 时的中心和非中心极大函数，并由此说明二者不相等。
\end{exercise}
\begin{solution}
中心区间半径 \(r\leq x-1\) 时积分为零；\(x-1<r<x\) 时平均为 \((r-x+1)/(2r)\)，随 \(r\) 增大；\(r\geq x\) 时平均为 \(1/(2r)\)，随 \(r\) 减小。所以 \(Mf(x)=1/(2x)\)。包含 \(x\) 的非中心区间若左端为 \(a\in[0,1]\)，平均至多 \((1-a)/(x-a)\leq1/x\)；左端小于零或右端超过 \(x\) 只增大分母。端点极限给出 \(M_uf(x)=1/x\)。
\end{solution}

% 自拟：组合两个分布尾部得到弱型乘积界，并检验强型结论的失败。
\Needspace{6\baselineskip}
\begin{exercise}\label{b2:ex:17-new-2}
设 \((X,\mathcal A,\mu)\) 为测度空间，\(0<p,q<\infty\)，并以 \(1/r=1/p+1/q\) 定义 \(r\)。对 \(f\in L^{p,\infty}(X)\)、\(g\in L^{q,\infty}(X)\)，证明
\[
 \|fg\|_{r,\infty}\leq 2^{1/r}\|f\|_{p,\infty}\cdot\|g\|_{q,\infty}.
\]
再在 \((0,1)\) 上构造这样的 \(f,g\)，使 \(fg\notin L^r(0,1)\)。不要求求出弱型乘积界的最佳常数。
\end{exercise}
\begin{solution}
记 \(A=\|f\|_{p,\infty}\)、\(B=\|g\|_{q,\infty}\)。若任一常数为零，相应函数几乎处处为零，结论成立。设 \(A,B>0\)。对任意 \(a,t>0\)，有
\[
 \{|fg|>t\}\subset\{|f|>a\}\cup\{|g|>t/a\},
 \quad d_{fg}(t)\leq A^pa^{-p}+B^qa^qt^{-q}.
\]
选择 \(a\) 使右边两项相等，即 \(a^{p+q}=A^pt^q/B^q\)。利用 \(r=pq/(p+q)\)，两项都等于 \((AB/t)^r\)，故 \(t\cdot d_{fg}(t)^{1/r}\leq2^{1/r}AB\)。对高度取上确界便得估计。

在 \((0,1)\) 上取 \(f(x)=x^{-1/p}\)、\(g(x)=x^{-1/q}\)。两者的弱拟范数均为一，而 \(|f(x)g(x)|^r=x^{-1}\)，其积分发散。因此两个弱型因子的乘积一般仍只能得到弱型控制。
\end{solution}

% 自拟：用统一弱型尾部与依测度收敛取得严格低于端点的范数收敛。
\Needspace{6\baselineskip}
\begin{exercise}\label{b2:ex:17-new-3}
设 \((X,\mathcal A,\mu)\) 为有限测度空间，\(0<p<\infty\)，可测函数列 \(f_n\) 满足 \(\sup_n\|f_n\|_{p,\infty}\leq A<\infty\)，且对每个 \(t>0\) 都有 \(\mu\{|f_n|>t\}\to0\)。证明对每个 \(0<r<p\)，都有 \(\|f_n\|_r\to0\)。在 \((0,1)\) 上给出反例，说明一般不能取 \(r=p\)。
\end{exercise}
\begin{solution}
记 \(m=\mu(X)\)，固定 \(R>0\)。层蛋糕公式把范数分成
\[
 \|f_n\|_r^r
 =r\int_0^R t^{r-1}d_{f_n}(t)\,\mathrm dt
  +r\int_R^\infty t^{r-1}d_{f_n}(t)\,\mathrm dt.
\]
第一段的被积函数逐点趋零，由可积函数 \(rm t^{r-1}\) 控制，因而趋零。第二段不超过
\[
 rA^p\int_R^\infty t^{r-p-1}\,\mathrm dt
 =\frac{rA^p}{p-r}R^{r-p}.
\]
先令 \(n\to\infty\)，再令 \(R\to\infty\)，得到结论。两种假设各有作用：依测度收敛处理有限高度，统一弱型界处理高值尾部。

反例取 \(f_n=n^{1/p}\cdot\mathbf1_{(0,1/n)}\)。它逐点趋零，弱 \(L^p\) 拟范数与强 \(L^p\) 范数都恒为一，而 \(\|f_n\|_r^r=n^{r/p-1}\to0\)。严格不等式 \(r<p\) 不能删除。
\end{solution}

% 自拟：精算幅度截断在弱Lp与L1中的不同误差。
\Needspace{6\baselineskip}
\begin{exercise}\label{b2:ex:17-new-4}
在 Lebesgue 测度空间 \((0,1)\) 上，设 \(f(x)=x^{-1/p}\)、\(1<p<\infty\)，并令 \(f_N(x)=\min\{f(x),N\}\)、\(N\geq1\)。求 \(f-f_N\) 的弱 \(L^p\) 拟范数；再求 \(\|f-f_N\|_1\) 并判断其极限。解释为什么这两个结论并不矛盾。
\end{exercise}
\begin{solution}
令 \(h_N=f-f_N=(f-N)_+\)。对 \(t>0\)，
\[
 d_{h_N}(t)=(N+t)^{-p},\quad
 \|h_N\|_{p,\infty}=\sup_{t>0}\frac{t}{N+t}=1.
\]
另一方面，层蛋糕公式给出
\[
 \|h_N\|_1=\int_0^\infty(N+t)^{-p}\,\mathrm dt
 =\frac{N^{1-p}}{p-1}\longrightarrow0.
\]
幅度截断消除了越来越少的 \(L^1\) 质量，但在任意高的水平上仍留下同一临界幂次尾部。弱拟范数取遍所有高度的上确界，因此不会随截断误差的积分一起趋零。此题也给出了上一题严格降低目标指数时的一个可算模型。
\end{solution}

% 自拟：将精细的极大函数尾部界用于LlogL输入的局部可积性。
\Needspace{6\baselineskip}
\begin{exercise}\label{b2:ex:17-new-5}
设 \(d\geq1\)，\(f:\mathbb R^d\to\mathbb C\) 可测且
\(\int_{\mathbb R^d}|f(x)|\log(e+|f(x)|)\,\mathrm dx<\infty\)。对任意有限 Lebesgue 测度集合 \(E\subset\mathbb R^d\)，证明中心极大函数满足
\[
 \int_E Mf\leq |E|+C_d\int_{\mathbb R^d}|f|\log(e+|f|).
\]
解释这一结论为什么没有推出 \(M:L^1\to L^1\)。
\end{exercise}
\begin{solution}
假设先保证 \(f\in L^1\)。对每个 \(t>0\)，把 \(f\) 分为大于 \(t/2\) 的部分和其余部分。后者的极大函数至多为 \(t/2\)，对前者用\cref{b2:thm:17-maximal-weak}，得
\[
 |\{Mf>t\}|\leq\frac{2\cdot5^d}{t}\int_{|f|>t/2}|f|.
\]
对 \(Mf\) 在 \(E\) 上用层蛋糕公式，高度 \(0<t<1\) 的积分至多为 \(|E|\)。其余部分用上式及非负积分换序，得到
\[
 \begin{aligned}
 \int_1^\infty |E\cap\{Mf>t\}|\,\mathrm dt
 &\leq2\cdot5^d\int_{|f|>1/2}|f(x)|
       \int_1^{2|f(x)|}\frac{\mathrm dt}{t}\,\mathrm dx\\
 &=2\cdot5^d\int_{|f|>1/2}|f|\log(2|f|).
 \end{aligned}
\]
当 \(s>1/2\) 时，\(\log(2s)\leq2\log(e+s)\)，故取 \(C_d=4\cdot5^d\) 即可。这里一方面对输入增加了对数可积性，另一方面只在有限测度集合上积分输出。即使 \(f\) 是单位球的示性函数，\(Mf\) 的全空间积分仍因远场尾部而发散。
\end{solution}

% 自拟：围绕本章已经建立的工具，改变条件或组合方法。
\Needspace{6\baselineskip}
\begin{exercise}\label{b2:ex:17-new-6}
设 \(d\geq1\)，\(K\in L^1(\mathbb R^d)\cap L^\infty(\mathbb R^d)\) 是复值核，\(Tf=K*f\)，Fourier 变换采用 \(e^{-2\pi\mathrm ix\cdot\xi}\) 的约定。用两个端点证明，对 \(1\leq p\leq2\) 及 \(f\in L^p(\mathbb R^d)\)，有
\[
 \|Tf\|_{p'}\leq\|K\|_\infty^{2/p-1}\cdot
 \|\widehat K\|_\infty^{2-2/p}\cdot\|f\|_p.
\]
其中 \(1/p+1/p'=1\)；\(K=0\) 时单独理解结论。
\end{exercise}
\begin{solution}
绝对积分给出 \(L^1\to L^\infty\) 常数 \(\|K\|_\infty\)。对 \(f\in L^1\cap L^2\)，卷积公式与 Plancherel 等式给出
\[
 \|K*f\|_2=\|\widehat K\,\widehat f\|_2
 \leq\|\widehat K\|_\infty\cdot\|f\|_2.
\]
由 \(L^2\) 稠密性得到平方端点；它与积分卷积在交空间上一致。对这两个兼容端点应用\cref{b2:thm:17-riesz-thorin}，源指数要求 \(1/p=1-\theta/2\)，故 \(\theta=2-2/p\)。目标倒数为 \(\theta/2=1/p'\)，恰得所需公式。一般输入的卷积还由 \(K\in L^1\) 的 Young 估计定义，它与上述延拓相同。若 \(\|\widehat K\|_\infty=0\)，Fourier 唯一性给出 \(K=0\)；若 \(\|K\|_\infty=0\)，也有同一结论，因而端点没有零的零次幂问题。
\end{solution}

% 自拟：围绕本章已经建立的工具，改变条件或组合方法。
\Needspace{6\baselineskip}
\begin{exercise}\label{b2:ex:17-new-7}
在测度空间 \((X,\mathcal A,\mu)\) 中，设 \(N\geq1\)，\(Q_1,\ldots,Q_N\) 是两两不交、有限正测度的可测集合，并记 \(U=\bigcup_jQ_j\)。对 \(f\in L^1(U;\mathbb C)\) 定义
\[
 Pf=\sum_{j=1}^N\left(\frac1{\mu(Q_j)}\int_{Q_j}f\,\mathrm d\mu\right)\cdot\mathbf1_{Q_j}.
\]
证明 \(P^2=P\)，且 \(P\) 在每个 \(L^p(U)\)、\(1\leq p\leq\infty\) 上的算子范数恰为一。进一步，若 \(f\in L^2(U)\)，证明 \(Pf\) 是所有在各 \(Q_j\) 上分别为常数的函数中唯一的最佳 \(L^2\) 逼近。
\end{exercise}
\begin{solution}
二次平均不改变常数，所以 \(P^2=P\)。积分三角不等式给出 \(\|Pf\|_1\leq\|f\|_1\)，每个平均不超过本性上确界又给出 \(\|Pf\|_\infty\leq\|f\|_\infty\)。在有限测度空间 \(U\) 上应用\cref{b2:thm:17-riesz-thorin}，所有中间指数的范数也至多为一。实函数可先把同一积分公式作用在复函数上，再限制回实函数。取非零函数 \(\mathbf1_U\)，它被 \(P\) 保持，故范数恰为一。

若 \(v\) 在各块上为常数，则每个块上 \(\int_{Q_j}(f-Pf)=0\)，从而
\(\int_U(f-Pf)\overline{(Pf-v)}\,\mathrm d\mu=0\)。展开平方模的积分便得
\[
 \|f-v\|_2^2=\|f-Pf\|_2^2+\|Pf-v\|_2^2.
\]
最小值在且仅在 \(v=Pf\) 几乎处处时取得。这同时说明为何平方误差下应取算术平均，而不必另猜每个块上的最优常数。
\end{solution}

% 自拟：用临界幂函数检验K泛函与Lorentz弱范数等价中的上界常数。
\Needspace{6\baselineskip}
\begin{exercise}\label{b2:ex:17-new-8}
在 Lebesgue 测度空间 \((0,1)\) 上，设 \(1<p<\infty\)、\(p'=p/(p-1)\)、\(f(x)=x^{-1/p}\)。计算 \(K(t,f;L^1,L^\infty)\)、\(t>0\)，并求
\[
 \sup_{t>0}t^{-1/p'}K(t,f;L^1,L^\infty).
\]
由此判断正文中 \(\|f\|_{1/p',\infty;K}\leq p'\|f\|_{p,\infty}\) 的系数 \(p'\) 能否普遍减小。进一步说明该函数为何不属于任何 \(L^{p,q}(0,1)\)、\(0<q<\infty\)。
\end{exercise}
\begin{solution}
此函数可积，重排在 \(0<s<1\) 时为 \(s^{-1/p}\)，在 \(s\geq1\) 时为零。由\cref{b2:prop:17-k-functional}，
\[
 K(t,f;L^1,L^\infty)=
 \begin{cases}
 p't^{1/p'},&0<t\leq1,\\
 p',&t>1.
 \end{cases}
\]
因此所求上确界为 \(p'\)，在每个 \(0<t\leq1\) 处取到，而 \(\|f\|_{p,\infty}=1\)。正文已经对任意输入证明了系数 \(p'\) 的上界，本例又达到该上界，所以这个等价估计的系数确实不能减小。

若 \(q<\infty\)，则 \(s^{1/p}f^*(s)=1\) 在 \((0,1)\) 上成立，Lorentz 量的 \(q\) 次幂包含 \(\int_0^1\mathrm ds/s=\infty\)。\(K\) 量也表现为所有小尺度上同一个常数：它的上确界有限，但对乘法测度 \(\mathrm dt/t\) 积分便发散。
\end{solution}
